\documentclass[11pt]{article}

% Basic packages (match series)
\usepackage[margin=1in]{geometry}
\usepackage{amsmath,amssymb,amsfonts}
\usepackage{bm}
\usepackage{graphicx}
\usepackage{hyperref}
\usepackage[numbers,sort&compress]{natbib}
\usepackage{authblk}

% Hyperref setup
\hypersetup{
  colorlinks=true,
  linkcolor=blue,
  citecolor=blue,
  urlcolor=blue
}

% Series macros (keep consistent with Papers I--VI)
\newcommand{\PN}{\mathrm{PN}}
\newcommand{\cS}{c_s}
\newcommand{\PhiP}{\Phi_{\mathrm{P}}}
\newcommand{\PhiL}{\Phi_{\mathrm{L}}}
\newcommand{\PhiTot}{\Phi}
\newcommand{\GM}{GM}
\newcommand{\ve}{\varepsilon}
\newcommand{\dd}{\mathrm{d}}

\title{Cosmological Sector from a Superfluid Defect Toy Model:\\
Finite-Slab Confinement, Flat Rotation Curves, and Optical Signatures}
\author{Trevor Norris}
\date{\today}

\begin{document}
\maketitle

\begin{abstract}
% TODO (1 paragraph):
% - What is the finite-slab mechanism?
% - What is proven: crossover 1/r^2 -> 1/(Hr), v_infty^2 = GM/H
% - What is added beyond dynamics: boundary sensitivity + lensing plateau
% - What is falsifiable: leakiness truncates plateau/flat window
\end{abstract}

\section{Introduction}
\label{sec:introduction}

Galaxy rotation curves and gravitational lensing are often interpreted as evidence that the
Newtonian/GR potential inferred from luminous matter is missing a large, extended component.
Standard explanations posit a nonluminous matter halo, while modified-gravity approaches seek an
infrared change in the effective force law.
In this paper we show that, within the brane--bulk superfluid defect toy model developed in the
preceding papers of this series, an infrared modification can arise \emph{geometrically} from a
finite transverse bulk extent, without introducing any additional ``dark'' source term.

The mechanism is the following.  Consider the scalar potential $\Phi$ sourced by a localized
mass/defect confined to the brane (taken as $z=0$), but assume the bulk is \emph{finite} in one
transverse direction, with walls at $z=\pm H$.
Imposing a no-flux (Neumann) boundary condition at the walls,
\begin{equation}
  \partial_z \Phi\big|_{z=\pm H}=0,
\end{equation}
changes the global Green's function while leaving the local near-source behavior unchanged.
Physically, field lines can spread spherically only up to radii comparable to the slab half-thickness
$H$; beyond that scale the dominant long-range behavior is controlled by the transverse zero mode,
which is effectively two-dimensional.

A central result of this work is that the brane-plane radial acceleration $F(r)$ transitions between
two controlled asymptotic regimes:
\begin{equation}
  F(r)\sim \frac{GM}{r^2}\quad (r\ll H),
  \qquad\qquad
  F(r)\sim \frac{GM}{H\,r}\quad (r\gg H).
  \label{eq:intro_force_asymptotes}
\end{equation}
The far-field $1/r$ scaling is the hallmark of dimensional reduction: the Neumann slab admits a true
transverse zero mode, so the long-range potential is logarithmic up to an additive constant (only
potential differences and gradients enter observables).
For circular motion in the brane plane,
\begin{equation}
  \frac{v^2(r)}{r}=F(r),
\end{equation}
Eq.~\eqref{eq:intro_force_asymptotes} implies a crossover from Keplerian decay $v(r)\propto r^{-1/2}$
to a flat asymptote
\begin{equation}
  v_\infty^2=\frac{GM}{H}.
  \label{eq:intro_vinf}
\end{equation}
Importantly, this flattening is a strict consequence of the slab Green's function; it is not a
phenomenological interpolation.

To make the claim numerically robust, we compute the same slab force in two mathematically distinct
representations: (i) a method-of-images construction (an infinite stack of mirror sources at
$z_n=2nH$) and (ii) a Neumann mode/Kaluza--Klein expansion, in which the force is a rapidly convergent
Bessel-$K$ series dominated at large $r$ by the zero mode.
Agreement between these two calculations, together with explicit near- and far-field normalization
checks, provides a nontrivial verification that the observed crossover is not a truncation artifact.
We also quantify the transition scale by a log-slope ``knee'' diagnostic, finding that the crossover
occurs at $r\sim H$ with an $\mathcal{O}(1)$ coefficient.

A second objective is to connect the slab regime to the optical sector developed earlier in the
series.  Using the same weak-field refractive-index mapping adopted in the optics paper, the bending
angle for an impact parameter $b$ can be written directly in terms of the brane acceleration profile:
\begin{equation}
  \Delta\theta(b)\simeq \frac{2}{c^2}\int_{-\infty}^{+\infty}
  F\!\left(\sqrt{b^2+z^2}\right)\,\frac{b}{\sqrt{b^2+z^2}}\,dz.
  \label{eq:intro_deflection_general}
\end{equation}
This makes the slab mechanism sharply testable.
In the Newtonian regime ($b\ll H$) one recovers the usual $\Delta\theta\propto 1/b$ scaling, while in
the deep slab regime ($b\gg H$) the $F\propto 1/r$ law implies a \emph{constant} deflection plateau,
\begin{equation}
  \Delta\theta(b)\to \Delta\theta_\infty=\frac{2\pi GM}{Hc^2}
  =\frac{2\pi v_\infty^2}{c^2},
  \qquad (b\gg H),
  \label{eq:intro_deflection_plateau}
\end{equation}
independent of impact parameter.
Equations \eqref{eq:intro_vinf} and \eqref{eq:intro_deflection_plateau} therefore provide a clean
lensing--rotation consistency relation that is specific to slab confinement.

Finally, we address robustness.
The $1/r$ force law is not generic: it relies on the existence of the Neumann zero mode.
If the walls are ``leaky'' (or effectively Dirichlet/absorbing), the zero mode is removed or lifted
to a small but nonzero transverse wavenumber, producing only a finite intermediate $1/r$ window
followed by eventual screening.
We present two practical parameterizations of leakage---a damped image stack and a Robin-boundary
mode spectrum---and show how both rotation-curve flattening and the lensing plateau degrade as
confinement weakens.

\paragraph{Organization.}
Section~\ref{sec:slab_bvp} formulates the slab boundary-value problem and presents the two equivalent
Green's-function representations (images and modes).
Section~\ref{sec:asymptotics_rotation} derives the asymptotic force laws and the implied rotation
curve behavior, including a quantitative definition of the transition (``knee'') scale.
Section~\ref{sec:leaky} studies boundary sensitivity and introduces leaky/Robin generalizations.
Section~\ref{sec:lensing} develops the optical predictions, emphasizing the deep-slab deflection
plateau and the lensing--rotation consistency relation.
Appendix~\ref{app:numerics} documents the numerical validation strategy (convergence studies and
cross-checks) used to support the main results.

\section{Inputs from Papers I--VI: dictionary and optical mapping}
\label{sec:inputs}

This paper is designed to be read as part of the 1PN series, but the slab-confinement mechanism
we study here is largely a statement about the \emph{Green's function} of the scalar (Newtonian)
sector and its optical consequences.  For self-containment we summarize, without re-derivation,
the minimal inputs we will reuse from Papers I--VI.

\subsection{Notation and dictionary}
We work with a brane located at $z=0$ embedded in a bulk with one transverse coordinate $z$.
In this paper the bulk is taken to be a finite slab,
\begin{equation}
  z \in [-H,H],
\end{equation}
with boundary conditions imposed at $z=\pm H$.
The in-brane radial coordinate is $r = \sqrt{x^2+y^2}$.

\paragraph{Line-of-sight coordinate in lensing.}
To avoid confusion with the bulk coordinate $z$, we denote the coordinate along an (approximately)
unperturbed light ray by $s$ in all deflection integrals (Paper II used $z$ for this role).
Thus, for impact parameter $b$,
\begin{equation}
  R(s) = \sqrt{b^2+s^2}
\end{equation}
will denote the 3D distance from the lens to a point on the fiducial straight path.

\paragraph{Mass parameter.}
Following the earlier papers, a localized defect/sink on the brane sources a Newtonian potential
with strength $\mu$.  When comparing to standard gravity we identify
\begin{equation}
  \mu \equiv GM,
\end{equation}
with $M$ the effective gravitational mass of the source.
We keep $\mu$ in intermediate equations for compactness and set $\mu=GM$ in the final GR-matching
expressions.

\paragraph{Wave/signal speed.}
The characteristic propagation speed in the superfluid vacuum is denoted $\cS$ (sound speed).
In the 1PN matching limit used throughout Papers I--VI, one sets $\cS=c$ when comparing to GR
orbital and optical observables.  We will follow the same convention here and explicitly retain
$c$ only in the optical formulas where it organizes the weak-field expansion.

\subsection{Scalar sector input: Poisson and lag potentials}
The scalar field sourced by brane defects is written as a sum
\begin{equation}
  \PhiTot = \PhiP + \PhiL,
\end{equation}
where $\PhiP$ is an instantaneous Poisson sector and $\PhiL$ is a finite-speed lag/retarded sector.
The series defines the strict Newtonian (0PN) limit as the regime in which the source is
time-independent and the lag sector can be neglected, so that $\PhiTot \to \PhiP$ and
\begin{equation}
  \nabla^2 \PhiP = 4\pi G\,\rho.
  \label{eq:inputs_poisson}
\end{equation}
For a point-like defect of strength $\mu$ located on the brane this yields
\begin{equation}
  \PhiP(r) = -\frac{\mu}{r}.
  \label{eq:inputs_point_potential}
\end{equation}

The present paper remains in this static/adiabatic regime: we are not using $\PhiL$ to generate the
infrared modification.  Instead, we keep the \emph{same} Poisson constraint
\eqref{eq:inputs_poisson} but change the global solution by changing the domain and boundary
conditions in the transverse direction.  Concretely, we will solve the Poisson/Laplace problem in a
finite slab $z\in[-H,H]$ with the source confined to $z=0$.  The near-source behavior remains
locally Newtonian, but the far field is controlled by the slab Green's function and its transverse
modes.

\paragraph{Force convention on the brane.}
For spherically (or cylindrically in-brane) symmetric configurations we define the \emph{inward}
radial acceleration magnitude on the brane by
\begin{equation}
  F(r) \equiv -\partial_r \Phi(r)\ge 0,
\end{equation}
so that circular orbits obey
\begin{equation}
  \frac{v^2(r)}{r} = F(r).
\end{equation}
In the Newtonian point-mass case \eqref{eq:inputs_point_potential}, $F(r)=\mu/r^2$.

\subsection{Optical mapping input (Paper II): refractive index and deflection}
Paper II treats lightlike excitations in the weak-field regime as rays propagating through an
inhomogeneous medium with refractive index
\begin{equation}
  N(\mathbf{x}) = \frac{c_0}{c_s(\mathbf{x})},
\end{equation}
where $c_s(\mathbf{x})$ is the local phase speed and $c_0$ is the asymptotic value.
In the 1PN matching limit one identifies $c_0\to c$, and lensing is computed to leading order in
$GM/(c^2 r)$.

For a point mass in the stiff-vacuum configuration used in Paper II, the result can be written as
\begin{equation}
  N(r) \simeq 1 + 2\,\frac{GM}{c^2 r},
  \qquad
  \ln N(r) \simeq 2\,\frac{GM}{c^2 r}.
  \label{eq:inputs_lnN_point}
\end{equation}
Since the Newtonian potential is $\Phi=-GM/r$, Eq.~\eqref{eq:inputs_lnN_point} is equivalently the
compact mapping
\begin{equation}
  \ln N(r) \simeq -\frac{2\Phi(r)}{c^2},
  \label{eq:inputs_lnN_general}
\end{equation}
valid to leading order in the weak-field expansion.
Only gradients of $\ln N$ enter ray deflection, so any additive constant in $\Phi$ (including the
gauge ambiguity of a logarithmic potential) has no optical effect at this order.

\paragraph{Deflection integral.}
In the thin-lens, small-angle approximation, the total bending angle accumulated by a ray with
impact parameter $b$ is (Paper II)
\begin{equation}
  \Delta\theta(b)
  \simeq \int_{-\infty}^{+\infty} \nabla_\perp \ln N\bigl(R(s)\bigr)\,\dd s,
  \label{eq:inputs_deflection_general}
\end{equation}
where $\nabla_\perp$ is the gradient transverse to the unperturbed ray direction.
For spherical symmetry,
\begin{equation}
  \left|\nabla_\perp \ln N\right|
  = \left|\frac{\dd \ln N}{\dd R}\right|\,\frac{b}{R}.
\end{equation}
Using \eqref{eq:inputs_lnN_general} and the force convention above, $\dd \Phi/\dd R = -F(R)$, this
can be rewritten in a form that will be central in the slab regime:
\begin{equation}
  \Delta\theta(b)
  \simeq \frac{2}{c^2}\int_{-\infty}^{+\infty}
  F\!\left(R(s)\right)\,\frac{b}{R(s)}\,\dd s,
  \qquad R(s)=\sqrt{b^2+s^2}.
  \label{eq:inputs_deflection_forceform}
\end{equation}
Equation~\eqref{eq:inputs_deflection_forceform} shows that, once the brane-plane acceleration law
$F(r)$ is determined (by any mechanism), its lensing consequences follow immediately with no
additional model freedom at leading order.

\paragraph{What we will and will not use from the earlier papers.}
The remainder of this work uses only:
(i) the static Poisson sector \eqref{eq:inputs_poisson} and the definition of the brane-plane force
$F(r)$, and
(ii) the optical mapping \eqref{eq:inputs_lnN_general}--\eqref{eq:inputs_deflection_forceform}.
We do not require the detailed 1PN kinetic-prefactor construction used to match perihelion
precession, nor the electromagnetic/vorticity sector.
The novelty here is that the same scalar Poisson constraint, when solved in a finite slab geometry,
induces an infrared crossover in $F(r)$ and a corresponding, sharply testable modification of the
lensing profile.

\section{Finite-slab scalar sector: geometry and boundary conditions}
\label{sec:slab_bvp}

This section defines the scalar boundary-value problem whose Green's function produces the infrared
crossover studied in the remainder of the paper.
We remain in the static (Poisson) sector summarized in
Section~\ref{sec:inputs} and modify only the global domain and boundary conditions in the transverse
direction.

\subsection{Slab geometry and Poisson problem}
\label{sec:slab_geometry}

We model the bulk as a finite slab in one transverse coordinate $z$,
\begin{equation}
  \mathcal{D} \;=\; \left\{ (x,y,z)\in\mathbb{R}^3 \,:\, z\in[-H,H] \right\},
\end{equation}
with the brane identified with the midplane $z=0$.
In the brane plane we use polar coordinates $(r,\varphi)$ with
\begin{equation}
  r=\sqrt{x^2+y^2}.
\end{equation}
All observables discussed here are brane-plane quantities evaluated at $z=0$.

\paragraph{Source localization on the brane.}
A localized defect/mass distribution on the brane is represented by a surface density
$\Sigma(x,y)$ (or $\Sigma(r)$ in axisymmetry) with support near the origin.
In the bulk Poisson sector, this corresponds to a 3D density
\begin{equation}
  \rho(x,y,z) = \Sigma(x,y)\,\delta(z).
\end{equation}
The static scalar potential $\Phi(r,z)$ then satisfies the Poisson equation in the slab,
\begin{equation}
  \nabla^2 \Phi(r,z) \;=\; 4\pi G\,\Sigma(r)\,\delta(z),
  \qquad (r,z)\in \mathbb{R}^2\times[-H,H],
  \label{eq:slab_poisson}
\end{equation}
where $\nabla^2=\partial_x^2+\partial_y^2+\partial_z^2$ is the ordinary Euclidean Laplacian.

\paragraph{Point-source normalization.}
For a point-like source of strength $\mu\equiv GM$ at $(r,z)=(0,0)$,
\begin{equation}
  \Sigma(r)=M\,\delta^{(2)}(\mathbf{r}) \quad \Longrightarrow \quad
  \nabla^2 \Phi = 4\pi \mu\,\delta^{(2)}(\mathbf{r})\,\delta(z).
\end{equation}
The normalization is chosen so that in the infinite-bulk limit $H\to\infty$ the solution reduces to
the standard Newtonian form,
\begin{equation}
  \Phi(\mathbf{x}) \to -\frac{\mu}{R},\qquad R=\sqrt{r^2+z^2}.
\end{equation}

\paragraph{Brane-plane force.}
The inward brane-plane acceleration magnitude is, by definition,
\begin{equation}
  F(r)\equiv -\partial_r \Phi(r,0)\ge 0,
  \label{eq:slab_force_def}
\end{equation}
and circular orbits satisfy $v^2(r)/r = F(r)$.

\subsection{Boundary conditions: sealed (Neumann) baseline and leaky generalizations}
\label{sec:slab_bc}

The only ingredient responsible for the infrared modification in this paper is the boundary
condition imposed at the slab walls $z=\pm H$.
In the ontology developed in \texttt{brane\_bulk\_ontology.tex}, the Poisson-sector scalar is an
effective potential for an irrotational bulk flow; a natural ``sealed'' condition is therefore a
no-flux wall, implemented as a Neumann boundary condition:
\begin{equation}
  \partial_z \Phi(r,z)\big|_{z=\pm H} = 0.
  \label{eq:bc_neumann}
\end{equation}
This is the baseline case we validate numerically and use for the main predictions.

To parameterize departures from perfect confinement we will also consider a one-parameter Robin
family,
\begin{equation}
  \partial_z \Phi(r,z)\big|_{z=\pm H} + \alpha\,\Phi(r,\pm H)=0,
  \label{eq:bc_robin}
\end{equation}
with $\alpha\ge 0$.
The limits of interest are:
\begin{equation}
  \alpha=0 \;\Rightarrow\; \text{Neumann (sealed)},\qquad
  \alpha\to\infty \;\Rightarrow\; \text{Dirichlet-like (strongly leaky/absorbing)}.
\end{equation}
We emphasize that \eqref{eq:bc_robin} is not introduced as a fit function; it is a controlled way to
track how the transverse mode spectrum changes as confinement weakens, and thereby how the infrared
force law changes.

\subsection{Why Neumann produces a $1/r$ far field: the zero-mode mechanism}
\label{sec:zero_mode_mechanism}

The qualitative reason the Neumann slab yields a $1/r$ force law on the brane is the existence of a
true transverse \emph{zero mode}.
To see this, consider separation of variables away from the source plane ($z\neq 0$), where
$\Phi$ is harmonic:
\begin{equation}
  \nabla^2 \Phi = 0 \qquad (z\neq 0).
\end{equation}
Writing a mode expansion
\begin{equation}
  \Phi(r,z) = \sum_{m=0}^\infty \phi_m(r)\,f_m(z),
  \label{eq:mode_ansatz}
\end{equation}
the transverse eigenfunctions satisfy
\begin{equation}
  -f_m''(z) = k_m^2\,f_m(z),
  \qquad z\in[-H,H],
  \label{eq:transverse_eigenproblem}
\end{equation}
with the chosen boundary condition at $z=\pm H$.

\paragraph{Neumann case.}
For \eqref{eq:bc_neumann}, the eigenproblem admits the constant function
\begin{equation}
  f_0(z) = \text{const},\qquad k_0=0,
\end{equation}
as an allowed mode.  This is the key: the long-distance behavior is then dominated by $m=0$ because
all higher modes have $k_m>0$ and are Yukawa/Bessel suppressed in the brane plane.
In the Green's-function language developed in the next section, the $k_0=0$ contribution produces a
logarithmic brane-plane potential at large $r$,
\begin{equation}
  \Phi(r,0) \sim -\frac{\mu}{H}\,\ln r \quad (r\gg H),
\end{equation}
and therefore a $1/r$ force,
\begin{equation}
  F(r)\equiv-\partial_r\Phi(r,0)\sim \frac{\mu}{H\,r}\quad (r\gg H).
\end{equation}
The additive constant ambiguity in $\ln r$ is physically irrelevant because observables depend on
gradients (forces) and on $\nabla_\perp \ln N$ in the optical sector.

\paragraph{Leaky case.}
For Robin boundaries \eqref{eq:bc_robin} with $\alpha>0$, the constant mode is no longer an exact
eigenfunction.  Equivalently, the Neumann zero mode is lifted to a small but nonzero transverse
wavenumber $k_0(\alpha)>0$.
As a result, the would-be logarithmic far field is eventually replaced by a screened behavior beyond
a length scale $\sim 1/k_0$.
This provides a clean robustness/falsifiability lever: persistent $1/r$ forces and plateau lensing
require sufficiently strong confinement, while imperfect confinement predicts an intermediate
``slab-like'' regime followed by suppression at the largest radii.

The remainder of the paper makes this mechanism quantitative by constructing the slab Green's
function in two equivalent representations (images and modes) and by explicitly tracking how the
infrared behavior changes as the boundary condition is deformed away from Neumann.

\section{Green's functions in the slab: images and KK/modes}
\label{sec:greens}

The infrared crossover derived in this paper is, at bottom, a statement about the Green's function of
the Laplacian in a finite slab with specified boundary conditions.  In the sealed (Neumann) case, the
same Green's function admits two complementary and numerically useful representations:
(i) a method-of-images construction and (ii) a transverse eigenmode (KK) expansion.  We present both,
because (a) each is convenient in a different regime and (b) their agreement provides a powerful
internal consistency check.

Throughout this section we treat a point source at $(r,z)=(0,0)$ with strength $\mu\equiv GM$.
Extended sources can be obtained by superposition once the point-source Green's function is known.

\subsection{Neumann Green's function as a boundary-value problem}
Let $G_N(\mathbf{r},z)$ denote the Neumann Green's function in the slab, defined by
\begin{align}
  \nabla^2 G_N(\mathbf{r},z) &= 4\pi\,\delta^{(2)}(\mathbf{r})\,\delta(z),
  \label{eq:GN_def}\\
  \partial_z G_N(\mathbf{r},z)\big|_{z=\pm H} &= 0,
  \label{eq:GN_neumannBC}
\end{align}
together with decay at $|\mathbf{r}|\to\infty$ in the sense appropriate for a slab (the potential
itself is defined only up to an additive constant, but gradients are finite and unique).
The physical potential is
\begin{equation}
  \Phi(\mathbf{r},z) = -\mu\,G_N(\mathbf{r},z),
\end{equation}
and the brane-plane force is computed from $F(r)=-\partial_r\Phi(r,0)$.

\subsection{Image-charge construction (sealed Neumann walls)}
\label{sec:images}

Neumann conditions at $z=\pm H$ are enforced by an even reflection of the source across each wall,
which is equivalent to a periodic array of identical sources with period $2H$ in the transverse
direction.  Concretely,
\begin{equation}
  G_N(r,z) \;=\; \sum_{n=-\infty}^{+\infty}\frac{1}{\sqrt{r^2+(z-2nH)^2}}
  \;+\; \text{(constant)}.
  \label{eq:GN_images}
\end{equation}
The additive constant reflects the fact that the slab potential is only defined up to a gauge shift;
it cancels in all observables we use (forces and deflections), and in practice we work directly with
gradients to avoid handling the constant explicitly.

On the brane ($z=0$), differentiating \eqref{eq:GN_images} gives a rapidly convergent and manifestly
positive expression for the inward radial acceleration magnitude:
\begin{equation}
  F_{\rm img}(r)
  \;=\; \mu\,\sum_{n=-\infty}^{+\infty}
  \frac{r}{\left(r^2+(2nH)^2\right)^{3/2}}.
  \label{eq:F_images}
\end{equation}
The $n=0$ term is the usual Newtonian $1/r^2$ force, and the additional images produce the infrared
modification.  For numerical evaluation, \eqref{eq:F_images} is naturally truncated at $|n|\le n_{\max}$
with an analytic tail estimate (Appendix~\ref{app:numerics}); the tail is particularly important when
$r\gg H$, where many images contribute comparably to produce the $1/(Hr)$ scaling.

\subsection{Mode / KK expansion}
\label{sec:modes}

An alternative representation is obtained by expanding in transverse eigenmodes of the Neumann
Sturm--Liouville problem on $[-H,H]$:
\begin{equation}
  f_m(z)=\cos(k_m z),\qquad k_m=\frac{m\pi}{H},\qquad m=0,1,2,\dots
  \label{eq:neumann_modes}
\end{equation}
where $m=0$ is the constant (zero) mode.
Writing the Green's function as a cosine series in $z$ and performing the 2D Fourier/Hankel
inversion in the brane plane yields the standard Yukawa/Bessel kernel $K_0$ for each $m\ge 1$.
On the brane ($z=0$), one convenient form for the potential is
\begin{equation}
  \Phi(r,0)
  \;=\; -\mu\left[
    -\frac{1}{H}\ln\!\left(\frac{r}{r_0}\right)
    + \frac{2}{H}\sum_{m=1}^{\infty}
      K_0\!\left(\frac{m\pi r}{H}\right)
  \right],
  \label{eq:Phi_modes}
\end{equation}
where $r_0$ is an arbitrary reference scale absorbing the additive constant (gauge choice).
Only $\partial_r\Phi$ will be used, so $r_0$ never enters final observables.

Differentiating \eqref{eq:Phi_modes} gives the mode-sum force law
\begin{equation}
  F_{\rm mod}(r)
  \;=\; -\partial_r\Phi(r,0)
  \;=\; \mu\left[
    \frac{1}{H\,r}
    + \frac{2\pi}{H^2}\sum_{m=1}^{\infty}
      m\,K_1\!\left(\frac{m\pi r}{H}\right)
  \right].
  \label{eq:F_modes}
\end{equation}
This representation makes the infrared physics transparent:
because $K_1(x)$ decays exponentially for large $x$, the sum is negligible for $r\gg H$ and the zero
mode gives the dominant far field,
\begin{equation}
  F(r)\;\xrightarrow{\,r\gg H\,}\; \frac{\mu}{H\,r}.
\end{equation}
Conversely, for $r\ll H$ many KK modes contribute and \eqref{eq:F_modes} reconstructs the local
Newtonian $1/r^2$ behavior.  Numerically, \eqref{eq:F_modes} converges extremely rapidly in the
far field (only a few $m$ are needed once $r\gtrsim H$), while \eqref{eq:F_images} can be preferable
very near the source where the $n=0$ term dominates.

\subsection{Equivalence and practical use}
\label{sec:equivalence}

Equations \eqref{eq:F_images} and \eqref{eq:F_modes} are two representations of the \emph{same}
Neumann slab Green's function; their equality can be understood as a Poisson-summation relation
connecting the periodic image stack to the discrete transverse spectrum.  In particular, one may
differentiate the corresponding identity for $G_N$ to obtain the brane-plane force identity
\begin{equation}
  \sum_{n=-\infty}^{+\infty}
  \frac{r}{\left(r^2+(2nH)^2\right)^{3/2}}
  \;=\;
  \frac{1}{H\,r}
  + \frac{2\pi}{H^2}\sum_{m=1}^{\infty}
  m\,K_1\!\left(\frac{m\pi r}{H}\right),
  \label{eq:images_modes_identity}
\end{equation}
which is precisely $F_{\rm img}(r)/\mu = F_{\rm mod}(r)/\mu$.

In the numerical validation section we exploit this equivalence as a stringent cross-check:
for a given $(H,\mu)$ and a wide range of $r$, the truncated image and truncated mode computations
agree to high precision once convergence controls (tail estimates and $m_{\max}$) are applied.
This ensures that the observed crossover from $F\sim \mu/r^2$ to $F\sim \mu/(Hr)$ is a property of
the slab Green's function itself rather than an artifact of truncation or representation choice.

\section{Force law and rotation curves from slab confinement}
\label{sec:asymptotics_rotation}

In this section we extract the dynamical consequences of the Neumann slab Green's function presented
in Section~\ref{sec:greens}.  The key object is the brane-plane inward acceleration magnitude
\begin{equation}
  F(r)\equiv -\partial_r\Phi(r,0)\ge 0,
\end{equation}
computed either from the image representation \eqref{eq:F_images} or the KK/mode representation
\eqref{eq:F_modes}.  We first establish the near- and far-field asymptotes (including their
normalizations) for a point source, then translate them into rotation-curve predictions, and finally
state the corresponding scaling for extended axisymmetric sources.

\subsection{Point mass: near-field and far-field asymptotes}
\label{sec:point_asymptotes}

For a point source of strength $\mu\equiv GM$ at $(r,z)=(0,0)$, the Neumann slab force can be written
in either of the equivalent forms:
\begin{align}
  F_{\rm img}(r)
  &= \mu\sum_{n=-\infty}^{+\infty}
  \frac{r}{\left(r^2+(2nH)^2\right)^{3/2}},
  \label{eq:dyn_F_images}\\[4pt]
  F_{\rm mod}(r)
  &= \mu\left[
    \frac{1}{H\,r}
    + \frac{2\pi}{H^2}\sum_{m=1}^{\infty}
      m\,K_1\!\left(\frac{m\pi r}{H}\right)
  \right].
  \label{eq:dyn_F_modes}
\end{align}
We denote the common value by $F(r)$.

\paragraph{Near field: $r\ll H$ (locally Newtonian).}
In the image form \eqref{eq:dyn_F_images}, the $n=0$ term is exactly Newtonian:
\begin{equation}
  F_{n=0}(r)=\frac{\mu}{r^2}.
\end{equation}
All other images satisfy $|2nH|\gg r$ and contribute a regular series in $r/H$.  Expanding the
$n\neq 0$ terms for $r\ll H$ gives
\begin{equation}
  \frac{r}{\left(r^2+(2nH)^2\right)^{3/2}}
  = \frac{r}{(2|n|H)^3}\left[1+\mathcal{O}\!\left(\frac{r^2}{n^2H^2}\right)\right],
\end{equation}
so that
\begin{equation}
  F(r)
  = \frac{\mu}{r^2}
    + \mu\left[\frac{\zeta(3)}{4}\,\frac{r}{H^3}
    + \mathcal{O}\!\left(\frac{r^3}{H^5}\right)\right],
  \qquad (r\ll H).
  \label{eq:dyn_nearfield_expansion}
\end{equation}
The leading correction is suppressed relative to $\mu/r^2$ by
$\sim (\zeta(3)/4)\,(r/H)^3$, so the force is locally indistinguishable from the usual $1/r^2$ law
well inside the slab thickness.  In particular,
\begin{equation}
  \lim_{r/H\to 0} r^2 F(r) = \mu,
\end{equation}
which fixes the near-field normalization unambiguously.

\paragraph{Far field: $r\gg H$ (slab regime).}
In the mode form \eqref{eq:dyn_F_modes}, the $m\ge 1$ contributions are exponentially suppressed for
$r\gg H$ because $K_1(x)\sim \sqrt{\pi/(2x)}\,e^{-x}$ as $x\to\infty$.  Therefore
\begin{equation}
  F(r)
  = \frac{\mu}{H\,r}\left[1+\mathcal{O}\!\left(e^{-\pi r/H}\right)\right],
  \qquad (r\gg H),
  \label{eq:dyn_farfield}
\end{equation}
so that
\begin{equation}
  \lim_{r/H\to\infty} H\,r\,F(r) = \mu.
\end{equation}
This establishes both the $1/r$ scaling and its normalization: the far-field force is
\begin{equation}
  F(r)\simeq \frac{GM}{H\,r},
\end{equation}
equivalently $\Phi(r,0)\simeq -(GM/H)\ln r$ up to an additive constant.

\paragraph{Crossover scale.}
The transition between \eqref{eq:dyn_nearfield_expansion} and \eqref{eq:dyn_farfield} occurs at
$r=\mathcal{O}(H)$, as expected from dimensional considerations: for $r\lesssim H$ field lines can
spread quasi-spherically, whereas for $r\gtrsim H$ the transverse direction is effectively saturated
and the dynamics are controlled by the Neumann zero mode.  In Section~\ref{sec:numerics} we quantify
this crossover with a log-slope knee diagnostic.

\subsection{Rotation curve consequence}
\label{sec:rotation_curves}

For circular motion in the brane plane,
\begin{equation}
  \frac{v^2(r)}{r} = F(r).
  \label{eq:dyn_circular_orbit}
\end{equation}
In the Newtonian regime $r\ll H$, Eqs.~\eqref{eq:dyn_nearfield_expansion} and
\eqref{eq:dyn_circular_orbit} give
\begin{equation}
  v^2(r) \simeq \frac{\mu}{r}
  \qquad (r\ll H),
\end{equation}
i.e.\ Keplerian decline $v(r)\propto r^{-1/2}$.

In the slab regime $r\gg H$, combining \eqref{eq:dyn_farfield} with \eqref{eq:dyn_circular_orbit}
gives a flat asymptote:
\begin{equation}
  v^2(r) \;\xrightarrow{\,r\gg H\,}\; v_\infty^2
  \equiv \frac{\mu}{H} = \frac{GM}{H}.
  \label{eq:dyn_vinf}
\end{equation}
Thus, in the sealed Neumann slab, asymptotically flat rotation curves are a direct consequence of
the slab Green's function and require no additional halo component.  The single geometric scale $H$
sets the plateau velocity for a given source strength $GM$.

\subsection{Extended sources and enclosed mass in the slab regime}
\label{sec:extended_sources}

Because the Poisson sector is linear, the slab potential for an extended brane source can be written
as a 2D convolution of the surface density with the slab Green's function restricted to the brane:
\begin{equation}
  \Phi(\mathbf{r},0) = -G\int \dd^2\mathbf{r}'\,
  \Sigma(\mathbf{r}')\,G_N\!\left(|\mathbf{r}-\mathbf{r}'|,0\right).
\end{equation}
In the deep slab regime (distances large compared with $H$), the Green's function reduces to the
effective 2D logarithmic form associated with the Neumann zero mode,
\begin{equation}
  G_N(\rho,0)\;\simeq\; -\frac{1}{H}\ln\!\left(\frac{\rho}{r_0}\right),
  \qquad (\rho\gg H),
\end{equation}
and therefore the brane-plane dynamics become approximately those of a two-dimensional potential
theory with effective coupling $G/H$.

For an axisymmetric surface density $\Sigma=\Sigma(r)$, define the enclosed mass
\begin{equation}
  M_{\rm enc}(r) \equiv 2\pi \int_0^{r} \Sigma(\rho)\,\rho\,\dd\rho.
\end{equation}
In the slab regime, the corresponding radial acceleration is approximately
\begin{equation}
  F(r)\;\simeq\; \frac{G\,M_{\rm enc}(r)}{H\,r},
  \qquad \text{(axisymmetric, slab regime)},
  \label{eq:dyn_extended_slab}
\end{equation}
so that
\begin{equation}
  v^2(r)\;\simeq\;\frac{G\,M_{\rm enc}(r)}{H}.
  \label{eq:dyn_extended_v2}
\end{equation}
If the mass distribution is effectively compact on the scale of the radii of interest (so that
$M_{\rm enc}(r)\to M_{\rm tot}$), \eqref{eq:dyn_extended_slab} reduces to the point-mass far field
\eqref{eq:dyn_farfield}.  If instead $M_{\rm enc}(r)$ continues to grow with $r$ over some range, then
$v^2(r)$ will track that growth while remaining insensitive to $r$ itself, a distinctive
``mass-tracking'' behavior of the slab regime.

Equation~\eqref{eq:dyn_extended_slab} should be understood as the leading large-$r/H$ behavior of the
Neumann slab Green's function for axisymmetric sources; corrections appear near the transition scale
$r\sim H$ and near sharp edges of the mass distribution.  These corrections can be computed by
numerical convolution with either representation of the Green's function (images or modes), but the
leading slab scaling \eqref{eq:dyn_extended_slab} is sufficient for the analytic predictions and the
lensing connection developed later.

Finally, we emphasize that the persistence of \eqref{eq:dyn_farfield}--\eqref{eq:dyn_extended_slab}
is contingent on the Neumann zero mode.  In Section~\ref{sec:leaky} we show how imperfect confinement
(removal or lifting of the zero mode) truncates the $1/r$ force regime and therefore limits the range
over which \eqref{eq:dyn_vinf} and \eqref{eq:dyn_extended_v2} apply.

\section{Numerical validation and convergence tests}
\label{sec:numerics}

The slab mechanism is an infrared modification produced by a global boundary-value problem, so a
numerical demonstration must establish more than a visually suggestive log--log plot.  In this
section we document three nontrivial validations:

\begin{enumerate}
  \item \textbf{Normalization checks} in both asymptotic regimes ($r\ll H$ and $r\gg H$),
  \item \textbf{Convergence control} under truncation of the image stack and of the KK/mode sum,
  \item \textbf{Representation cross-check}: images and modes agree to numerical precision over a wide
        dynamic range.
\end{enumerate}

All results reported below are produced by the accompanying Mathematica validation script
(\texttt{dark\_matter\_slab\_test\_revised.wl}) using a point source on the brane and the Neumann slab
formulas \eqref{eq:dyn_F_images}--\eqref{eq:dyn_F_modes}.
Unless stated otherwise, we use $H=10$ and $\mu=GM=1$.

\subsection{Critical normalization checks}
\label{sec:numerics_norm}

The point-mass slab force has two sharp normalization limits:
\begin{equation}
  r^2 F(r)\to \mu \quad (r\ll H),
  \qquad
  Hr\,F(r)\to \mu \quad (r\gg H).
  \label{eq:numerics_norm_targets}
\end{equation}
These limits are stringent because they test both the scaling and the overall coefficient.

Using both the image representation and the mode representation, we evaluate the normalized products
at representative near- and far-field radii.  With $H=10$ and $\mu=1$ we obtain:
\begin{align}
  r=0.5\ (r\ll H): \quad & r^2 F_{\rm img}(r)=1.000037534,\qquad r^2 F_{\rm mod}(r)=1.000037534,\\
  r=1000\ (r\gg H): \quad & HrF_{\rm img}(r)=1.000000000,\qquad HrF_{\rm mod}(r)=1.000000000,
\end{align}
confirming the expected asymptotes \eqref{eq:dyn_nearfield_expansion} and \eqref{eq:dyn_farfield} with
the correct normalization.

\subsection{Convergence study: truncation control}
\label{sec:numerics_convergence}

Both numerical representations involve truncation.  In the image method
\eqref{eq:dyn_F_images}, we sum images over $|n|\le n_{\max}$ and add a controlled tail estimate for
$|n|>n_{\max}$; in the mode method \eqref{eq:dyn_F_modes}, we truncate the KK series at
$m\le m_{\max}$.
To quantify convergence we compare truncated results to a high-accuracy reference computed with a
larger cutoff (and, in the image case, an explicit analytic tail).

\paragraph{Image representation.}
Define $F_{\rm img}^{(n_{\max})}(r)$ as the truncated image sum with tail correction and let
$F_{\rm ref}(r)$ be a stricter reference (larger cutoff and tighter tail tolerance).  We report the
relative error
\begin{equation}
  \varepsilon_{\rm img}(r;n_{\max}) \equiv
  \left|\frac{F_{\rm img}^{(n_{\max})}(r)-F_{\rm ref}(r)}{F_{\rm ref}(r)}\right|.
\end{equation}
Representative values (full table in Appendix~\ref{app:numerics}) show rapid improvement with
$n_{\max}$, with large-$r$ cases requiring more images as expected:
\begin{center}
\begin{tabular}{c|ccccc}
\hline
$r$ & $n_{\max}=200$ & $500$ & $1000$ & $2000$ & $5000$ \\
\hline
$1$     & $1.9\times 10^{-14}$ & $4.4\times 10^{-16}$ & $\sim 0$ & $\sim 0$ & $\sim 0$ \\
$100$   & $1.9\times 10^{-9}$  & $5.0\times 10^{-11}$ & $3.1\times 10^{-12}$ & $1.9\times 10^{-13}$ & $5.4\times 10^{-15}$ \\
$1000$  & $1.7\times 10^{-7}$  & $4.9\times 10^{-9}$  & $3.1\times 10^{-10}$ & $1.9\times 10^{-11}$ & $5.0\times 10^{-13}$ \\
\hline
\end{tabular}
\end{center}
The pattern matches the physics: deep in the slab regime ($r\gg H$), many images contribute
coherently to build the $1/(Hr)$ law, so the raw image sum requires a larger cutoff; the tail
estimate restores accuracy efficiently.

\paragraph{Mode representation.}
Define $F_{\rm mod}^{(m_{\max})}(r)$ as the truncated KK sum and measure
\begin{equation}
  \varepsilon_{\rm mod}(r;m_{\max}) \equiv
  \left|\frac{F_{\rm mod}^{(m_{\max})}(r)-F_{\rm mod}^{(m_{\rm ref})}(r)}{F_{\rm mod}^{(m_{\rm ref})}(r)}\right|,
\end{equation}
with a strict reference cutoff $m_{\rm ref}$.
Because $K_1(x)$ decays exponentially for large $x$, convergence is typically extremely fast once
$r\gtrsim H$.  In our runs, $m_{\max}\simeq 100$ already yields relative errors consistent with
machine precision for essentially all $r\ge 1$; even $m_{\max}=50$ is sufficient over most of the
range (see Appendix~\ref{app:numerics} for the full table).

These convergence tests support a practical division of labor: images are conceptually transparent
and efficient very near the source, while KK/modes are highly efficient in the far field where only
the zero mode and a few low-$m$ corrections matter.

\subsection{Cross-check: images vs modes}
\label{sec:numerics_crosscheck}

The most stringent validation is agreement between the two independent representations when each is
evaluated at its converged cutoff.  For a representative selection of radii, we compare
$F_{\rm img}(r)$ and $F_{\rm mod}(r)$ and report the relative difference
\begin{equation}
  \delta(r)\equiv
  \left|\frac{F_{\rm mod}(r)-F_{\rm img}(r)}{F_{\rm img}(r)}\right|.
\end{equation}
With $H=10$ and $\mu=1$, we find:
\begin{center}
\begin{tabular}{c|ccc}
\hline
$r$ & $F_{\rm img}(r)$ & $F_{\rm mod}(r)$ & $\delta(r)$ \\
\hline
$1$    & $1.0002995450516297$ & $1.0002995450516299$ & $2.2\times 10^{-16}$ \\
$10$   & $0.012261662448154562$ & $0.012261662448154548$ & $1.1\times 10^{-15}$ \\
$100$  & $1.000000000000329\times 10^{-3}$ & $1.000000000000323\times 10^{-3}$ & $6.1\times 10^{-15}$ \\
$1000$ & $1.0000000000004988\times 10^{-4}$ & $1.0000000000000000\times 10^{-4}$ & $5.0\times 10^{-13}$ \\
\hline
\end{tabular}
\end{center}
Agreement at the $10^{-15}$--$10^{-13}$ level over three decades in radius indicates that the
observed crossover is not a representation artifact.  The slightly larger relative discrepancy at
$r=1000$ is consistent with the image method being intrinsically more demanding in the deep slab
regime unless a very large $n_{\max}$ (plus tail) is used.

\subsection{Knee (transition) scale estimate}
\label{sec:numerics_knee}

To quantify the crossover radius without relying on visual inspection of plots, we use a log-slope
diagnostic.  Define
\begin{equation}
  \gamma(r)\equiv \frac{\dd \ln F}{\dd \ln r}.
\end{equation}
In the Newtonian regime, $\gamma\to -2$; in the slab regime, $\gamma\to -1$.
We define the ``knee'' radius $r_{\rm knee}$ as the location where $\gamma(r_{\rm knee})=-3/2$.
For the representative run with $H=10$, the script finds
\begin{equation}
  r_{\rm knee}\simeq 10.50,
\end{equation}
consistent with the expectation that the transition occurs at $r=\mathcal{O}(H)$ with an order-unity
coefficient.

\paragraph{Summary.}
The combined evidence---asymptotic normalization, controlled convergence, and images--modes
agreement---supports treating the slab crossover
\begin{equation}
  F(r)\sim \frac{\mu}{r^2}\ (r\ll H)
  \quad\longrightarrow\quad
  F(r)\sim \frac{\mu}{H\,r}\ (r\gg H)
\end{equation}
as a robust consequence of the Neumann slab Green's function, suitable for use in the prediction
sections that follow.

\section{Boundary sensitivity: leaky slab models}
\label{sec:leaky}

The Neumann slab result is a \emph{conditional} geometric statement:
a persistent far-field $F(r)\sim \mu/(Hr)$ requires a transverse zero mode.
In the toy ontology, this corresponds to a bulk that is effectively ``sealed'' to the Poisson-sector
flow (no normal flux through $z=\pm H$).  If the boundary is imperfect, the infrared behavior is
modified in a controlled and testable way.

The goal of this section is not to introduce new free functions, but to provide two simple,
implementable parameterizations of ``leakiness'' and to classify the resulting phenomenology.
Both models reduce to the Neumann slab in an appropriate limit and both produce a clear,
falsifiable prediction: \emph{the slab-like $1/r$ regime (and its lensing plateau) becomes an
intermediate window rather than an asymptote.}

\subsection{What is being tested: the fate of the zero mode}
\label{sec:leaky_zeromode}

Section~\ref{sec:zero_mode_mechanism} showed that Neumann confinement admits a true transverse
zero mode ($k_0=0$).  In the KK representation \eqref{eq:dyn_F_modes}, this is the origin of the
explicit $1/(Hr)$ term.  Any departure from perfect confinement does one of two things:

\begin{itemize}
\item \textbf{Reduces the coherence of the image stack} (``partial reflection''): the effective
      number of images contributing to the far field becomes finite.  One expects an intermediate
      slab-like window followed by a return toward a more Newtonian-like falloff at the largest
      radii.
\item \textbf{Lifts the zero mode} (``absorptive/leaky walls''): the would-be $k_0=0$ mode becomes a
      small but nonzero $k_0>0$, leading to eventual \emph{screening} at $r\gtrsim 1/k_0$.
\end{itemize}

We now implement these two possibilities as Model~A and Model~B.

\subsection{Model A: weighted images (partial reflection)}
\label{sec:leaky_weighted_images}

A fast and intuitive way to model imperfect confinement is to modify the Neumann image stack by an
amplitude factor that decays with image index:
\begin{equation}
  F_{w}(r)
  \;=\;
  \mu \sum_{n=-\infty}^{+\infty} w^{|n|}
  \frac{r}{\left(r^2+(2nH)^2\right)^{3/2}},
  \qquad 0<w\le 1.
  \label{eq:leaky_weighted_images_force}
\end{equation}
Here $w$ acts as a \emph{reflectivity} parameter:
\begin{equation}
  w=1 \Rightarrow \text{perfect Neumann slab (sealed)},
  \qquad
  w<1 \Rightarrow \text{partial leakage (finite effective image depth)}.
\end{equation}
Because the weights decay geometrically, the dominant contribution comes from
$|n|\lesssim n_{\rm eff}$ with
\begin{equation}
  n_{\rm eff}\sim \frac{1}{1-w},
  \label{eq:leaky_neff}
\end{equation}
so that the slab-like behavior is expected only over radii for which a large fraction of these
images contribute coherently.

\paragraph{Qualitative regimes.}
Model~A naturally produces three regimes:

\begin{enumerate}
\item \textbf{Near field} ($r\ll H$): the $n=0$ term dominates and
\begin{equation}
  F_w(r) = \frac{\mu}{r^2}\left[1+\mathcal{O}\!\left(\frac{r^3}{H^3}\right)\right].
\end{equation}

\item \textbf{Intermediate slab window} ($H\ll r\ll 2n_{\rm eff}H$): many weighted images contribute.
The force approaches a slab-like $1/r$ behavior over a finite dynamic range:
\begin{equation}
  F_w(r)\approx \frac{\mu}{H\,r}\times \mathcal{A}(w),
\end{equation}
where $\mathcal{A}(w)$ is an order-unity amplitude that approaches $1$ as $w\to 1$.

\item \textbf{Ultra-far field} ($r\gg 2n_{\rm eff}H$): the effective image depth is finite and the
sum behaves increasingly like a finite collection of sources, so the force trends back toward a
more Newtonian-like falloff (up to an overall renormalization of $\mu$).
\end{enumerate}

\paragraph{Diagnostics.}
To characterize the ``flat window'' we use the log-slope
\begin{equation}
  \gamma(r) \equiv \frac{\dd\ln F}{\dd\ln r},
\end{equation}
and define a slab-like interval by the condition $|\gamma(r)+1|<\epsilon$ for a fixed tolerance
(e.g.\ $\epsilon=0.1$).  The window length grows rapidly as $w\to 1$ through the scaling
\eqref{eq:leaky_neff}.  This makes $w$ directly interpretable as a confinement diagnostic rather than
a mere fit parameter.

\paragraph{Rotation curves.}
Within the slab-like window, circular speeds satisfy
\begin{equation}
  v^2(r)\simeq rF_w(r)\approx \frac{\mu}{H}\,\mathcal{A}(w),
\end{equation}
so rotation curves are approximately flat.  Outside the window, the curve either declines or
transitions toward a different asymptotic depending on $w$.

\subsection{Model B: Robin boundary condition and lifted zero mode}
\label{sec:leaky_robin}

A more formal description of leakage is to replace the Neumann wall with a Robin condition
\eqref{eq:bc_robin},
\begin{equation}
  \partial_z \Phi + \alpha \Phi = 0 \qquad \text{at } z=\pm H,
\end{equation}
with $\alpha>0$ encoding absorption/leakiness.

\paragraph{Eigenvalue condition.}
For even modes (appropriate for a brane source at $z=0$), the transverse eigenfunctions are of the
form $f(z)=\cos(k z)$, and the Robin condition yields the root condition
\begin{equation}
  -k\sin(kH)+\alpha\cos(kH)=0
  \quad\Longleftrightarrow\quad
  \tan(kH)=\frac{\alpha}{k}.
  \label{eq:robin_root_condition}
\end{equation}
For $\alpha=0$, Eq.~\eqref{eq:robin_root_condition} admits $k_0=0$ (the Neumann zero mode).
For $\alpha>0$, the lowest root becomes small but nonzero ($k_0>0$), i.e.\ the zero mode is lifted.

\paragraph{Effective screened far field.}
In the KK representation, the $k_0=0$ term that produced the logarithmic potential is replaced by a
screened two-dimensional kernel.  The leading far-field contribution takes the form
\begin{equation}
  \Phi_0(r,0)\propto -\frac{\mu}{H}\,K_0(k_0 r),
  \qquad
  F_0(r)\propto \frac{\mu}{H}\,k_0 K_1(k_0 r),
  \label{eq:robin_screened_force}
\end{equation}
so that for $k_0 r\gg 1$,
\begin{equation}
  F_0(r) \sim \frac{\mu}{H}\sqrt{\frac{\pi k_0}{2r}}\,e^{-k_0 r}.
  \label{eq:robin_screened_asymptote}
\end{equation}
Thus, the model predicts a slab-like regime only for $r\lesssim k_0^{-1}$, followed by exponential
suppression at the largest radii.

\paragraph{Small-leakage scaling.}
For weak leakage ($\alpha H\ll 1$), the smallest root of \eqref{eq:robin_root_condition} is small and
a useful scaling estimate is
\begin{equation}
  k_0 \sim \sqrt{\frac{\alpha}{H}},
  \label{eq:k0_scaling}
\end{equation}
so the screening length grows as $\ell_{\rm scr}\sim 1/k_0\sim \sqrt{H/\alpha}$.

\paragraph{Rotation curves.}
Within the intermediate range $H\ll r\ll \ell_{\rm scr}$, the force is approximately $F\sim \mu/(Hr)$
and rotation curves are quasi-flat.  For $r\gtrsim \ell_{\rm scr}$, the force becomes too small to
support a constant-speed asymptote, and $v(r)$ declines.

\subsection{Robustness taxonomy and falsifiable predictions}
\label{sec:leaky_taxonomy}

The two leaky models lead to a common qualitative conclusion: the Neumann slab asymptote is
\emph{diagnostic} of confinement.

\begin{itemize}
\item \textbf{Sealed (Neumann):} persistent $F\sim \mu/(Hr)$ at $r\gg H$, implying
      $v_\infty^2=\mu/H$ and (Section~\ref{sec:lensing}) a true large-$b$ lensing plateau.

\item \textbf{Partially reflecting (weighted images):} an intermediate slab-like window exists for
      $H\ll r\ll 2n_{\rm eff}H$, with its extent controlled by $w$ (or equivalently $n_{\rm eff}$).
      Beyond this window the force trends away from $1/r$.

\item \textbf{Absorbing/leaky (Robin, lifted zero mode):} the slab-like window ends at a screening
      scale $\ell_{\rm scr}\sim 1/k_0$, beyond which $F(r)$ is exponentially suppressed.
\end{itemize}

In both cases, the key observable is the \emph{extent} of the slab-like window in dynamics and optics.
In Section~\ref{sec:lensing} we show that the same confinement physics controlling the flattening of
rotation curves also controls the onset and persistence of the lensing plateau.  Joint rotation+%
lensing measurements therefore constrain not only the geometric scale $H$ but also (in the leaky
extensions) the leakage parameters $w$ or $k_0$ that quantify how ``sealed'' the bulk is.

\section{Optical predictions: lensing from the slab force}
\label{sec:lensing}

A central advantage of the slab mechanism is that it modifies the brane-plane acceleration law
$F(r)$ in a way that propagates directly into the optical sector with essentially no additional model
freedom at leading order.
This section derives the bending angle $\Delta\theta(b)$ as a functional of $F(r)$, extracts the
asymptotic lensing scalings in the Newtonian and slab regimes, and identifies a sharp signature of
slab confinement: a \emph{large-impact-parameter deflection plateau}.
We then summarize how imperfect confinement (Section~\ref{sec:leaky}) truncates this plateau.

\subsection{Deflection written directly in terms of $F(r)$}
\label{sec:lensing_force_form}

From the optical mapping summarized in Section~\ref{sec:inputs}, the weak-field refractive index
obeys
\begin{equation}
  \ln N \simeq -\frac{2\Phi}{c^2},
\end{equation}
and the small-angle deflection for a ray with impact parameter $b$ is
\begin{equation}
  \Delta\theta(b)\simeq \int_{-\infty}^{+\infty}\nabla_\perp \ln N\bigl(R(s)\bigr)\,\dd s,
  \qquad R(s)=\sqrt{b^2+s^2}.
\end{equation}
For spherical symmetry, $\nabla_\perp \ln N = (\dd \ln N/\dd R)\,(b/R)$.
Using $\dd\Phi/\dd R=-F(R)$ (our convention that $F\ge 0$ is the inward acceleration magnitude),
we obtain the key identity:
\begin{equation}
  \boxed{
  \Delta\theta(b)
  \;\simeq\;
  \frac{2}{c^2}\int_{-\infty}^{+\infty}
  F\!\left(R(s)\right)\,\frac{b}{R(s)}\,\dd s
  }
  \label{eq:lensing_force_identity}
\end{equation}
This formula is particularly well-suited to the present paper because $F(r)$ is the quantity already
validated by the slab Green's function computations.  Note also that any additive constant in $\Phi$
(including the gauge ambiguity of a logarithmic far-field potential) drops out of
\eqref{eq:lensing_force_identity}; only gradients matter.

\paragraph{Numerically stable finite-range form.}
For direct computation, it is convenient to map the infinite line-of-sight integral to a finite
interval using $s=b\tan u$, for which $u\in(-\pi/2,\pi/2)$ and $R=b\sec u$.  Then
\begin{equation}
  \dd s = b\,\sec^2 u\,\dd u,
  \qquad
  \frac{b}{R} = \cos u,
\end{equation}
and \eqref{eq:lensing_force_identity} becomes
\begin{equation}
  \boxed{
  \Delta\theta(b)
  \;=\;
  \frac{2b}{c^2}\int_{-\pi/2}^{+\pi/2} F\!\left(b\sec u\right)\,\sec u\,\dd u
  }
  \label{eq:lensing_uform}
\end{equation}
The integrand in \eqref{eq:lensing_uform} is smooth for all $u\in(-\pi/2,\pi/2)$ for any physically
regularized force law, and the finite limits make the computation robust and fast.

\subsection{Asymptotes: $1/b$ bending to a constant plateau}
\label{sec:lensing_asymptotes}

We now evaluate \eqref{eq:lensing_force_identity} in the two asymptotic regimes established in
Section~\ref{sec:asymptotics_rotation} for the sealed Neumann slab.

\paragraph{Near field: $b\ll H$ (Newtonian bending).}
When the ray probes only radii $R\ll H$ along its trajectory, the force is locally Newtonian:
$F(R)\simeq \mu/R^2$.
Then
\begin{align}
  \Delta\theta(b)
  &\simeq \frac{2}{c^2}\int_{-\infty}^{+\infty}
  \frac{\mu}{R^2}\,\frac{b}{R}\,\dd s
  = \frac{2\mu}{c^2}\int_{-\infty}^{+\infty}
  \frac{b}{(b^2+s^2)^{3/2}}\,\dd s
  = \frac{4\mu}{b\,c^2},
\end{align}
i.e.
\begin{equation}
  \Delta\theta(b)\;\xrightarrow{\,b\ll H\,}\; \frac{4GM}{b\,c^2}.
  \label{eq:lensing_newtonian}
\end{equation}
This reproduces the standard $1/b$ bending scaling in the regime where the slab is irrelevant.

\paragraph{Deep slab: $b\gg H$ (deflection plateau).}
In the slab regime, the brane-plane force obeys
\begin{equation}
  F(R)\simeq \frac{\mu}{H\,R},
  \qquad (R\gg H).
\end{equation}
Substituting into \eqref{eq:lensing_force_identity} yields
\begin{align}
  \Delta\theta(b)
  &\simeq \frac{2}{c^2}\int_{-\infty}^{+\infty}
  \frac{\mu}{H\,R}\,\frac{b}{R}\,\dd s
  = \frac{2\mu}{Hc^2}\int_{-\infty}^{+\infty}
  \frac{b}{b^2+s^2}\,\dd s
  = \frac{2\mu}{Hc^2}\,\pi.
\end{align}
Therefore the bending angle tends to a constant independent of $b$:
\begin{equation}
  \boxed{
  \Delta\theta(b)\;\xrightarrow{\,b\gg H\,}\;
  \Delta\theta_\infty
  \equiv \frac{2\pi\,\mu}{Hc^2}
  = \frac{2\pi GM}{Hc^2}
  }
  \label{eq:lensing_plateau}
\end{equation}
Combining \eqref{eq:lensing_plateau} with the flat rotation curve asymptote
$v_\infty^2=\mu/H$ from Section~\ref{sec:asymptotics_rotation} gives an especially compact
consistency relation:
\begin{equation}
  \boxed{
  \Delta\theta_\infty = \frac{2\pi\,v_\infty^2}{c^2}
  }
  \label{eq:lensing_rotation_relation}
\end{equation}
This is a sharp signature of slab confinement: in the same regime where $v(r)$ is asymptotically
flat, lensing bends toward an impact-parameter independent plateau.

\paragraph{Slope diagnostic.}
Defining the log-slope
\begin{equation}
  \beta(b)\equiv \frac{\dd\ln \Delta\theta}{\dd\ln b},
\end{equation}
Eqs.~\eqref{eq:lensing_newtonian} and \eqref{eq:lensing_plateau} imply a transition
\begin{equation}
  \beta(b)\to -1 \quad (b\ll H),
  \qquad
  \beta(b)\to 0 \quad (b\gg H).
\end{equation}
Thus, in addition to $\Delta\theta(b)$ itself, the slope $\beta(b)$ provides a clean and
scale-localized diagnostic of the crossover.

\subsection{Leakiness dependence of the lensing plateau}
\label{sec:lensing_leaky}

Section~\ref{sec:leaky} showed that imperfect confinement alters the fate of the Neumann zero mode.
Because the deflection \eqref{eq:lensing_force_identity} is a line-of-sight average of $F(R)$, the
same confinement physics controlling the rotation-curve flattening controls the onset and
persistence of the lensing plateau.

\paragraph{Model A (weighted images): plateau becomes a finite window.}
For the weighted-image force $F_w(r)$ in \eqref{eq:leaky_weighted_images_force}, the force can
approximate $F\propto 1/r$ only over a finite radial interval whose extent grows as $w\to 1$.
Consequently, $\Delta\theta(b)$ exhibits an \emph{approximate} plateau only for impact parameters
$b$ that predominantly sample this window.  As $b$ increases beyond the window, the deflection
departs from \eqref{eq:lensing_plateau} and begins to decrease.

\paragraph{Model B (Robin / lifted zero mode): plateau ends at a screening scale.}
For Robin leakage, the far field is controlled by a lifted lowest mode with $k_0>0$, producing a
screened contribution of the form \eqref{eq:robin_screened_force}.  In this case a slab-like window
(and an approximately constant $\Delta\theta$) exists only for
\begin{equation}
  H \ll b \ll \ell_{\rm scr}\equiv k_0^{-1},
\end{equation}
while for $b\gtrsim \ell_{\rm scr}$ the bending is suppressed in concert with the exponential decay
of the force.  Thus, measuring the \emph{extent} of any observed plateau provides a direct constraint
on confinement quality (via $w$ or $k_0$), while the plateau height constrains $H$ (or equivalently
$v_\infty$ through \eqref{eq:lensing_rotation_relation}).

\paragraph{Practical computational program.}
Given any converged numerical representation of $F(r)$ (images, modes, or leaky variants),
Eq.~\eqref{eq:lensing_uform} can be evaluated efficiently for a set of $b$ values to generate:
(i) $\Delta\theta(b)$ and its approach to the plateau,
(ii) the slope $\beta(b)$ as a crossover diagnostic, and
(iii) direct comparisons between sealed and leaky boundaries.
These curves are the primary falsifiable optical predictions of the slab-confinement mechanism.

\section{Cosmological embedding and interpretation}
\label{sec:cosmo}

The preceding sections established a controlled dynamical and optical statement:
solving the static scalar (Poisson) sector in a finite slab with Neumann walls produces an infrared
crossover on the brane from $F\propto r^{-2}$ to $F\propto (Hr)^{-1}$, implying flat rotation curves
and a large-$b$ lensing plateau.  This section addresses how the geometric parameter $H$ should be
interpreted in the wider ontology, what can be inferred from observations, and how the same picture
can be embedded into a cosmological-scale narrative without mixing assumptions into the core slab
result.

\subsection{Interpreting $H$ as a control parameter}
\label{sec:cosmo_H}

In this paper $H$ is defined operationally: it is the half-thickness of the region in the transverse
bulk direction that is dynamically accessible to the Poisson-sector scalar field sourced by brane
defects.  Equivalently, $H$ is the transverse confinement scale that sets both
(i) the \emph{knee} radius of the force-law crossover and
(ii) the normalization of the slab-regime coupling $G/H$.

Two distinct interpretations are possible (and observationally distinguishable):

\paragraph{(i) Global thickness: $H$ is approximately universal.}
If the brane resides in a bulk cavity of roughly fixed thickness, then $H$ is a global parameter of
the vacuum.
In this case, for sufficiently isolated/compact sources, the deep-slab asymptotes become
\begin{equation}
  v_\infty^2=\frac{GM}{H},
  \qquad
  \Delta\theta_\infty=\frac{2\pi GM}{Hc^2},
\end{equation}
so $H$ can be inferred directly from any system for which $M$ is independently known.
A universal $H$ would then be a strong global signature of the slab mechanism.

\paragraph{(ii) Effective thickness: $H$ varies with environment or scale.}
Alternatively, $H$ may be an \emph{effective} confinement scale that depends on location, epoch, or
the local configuration of the brane--bulk system (e.g.\ variations in bulk pressure, brane tension,
or the density of active throats).
In this case, the slab predictions remain valid \emph{locally} with $H\to H_{\rm eff}$ for a given
system, but one should not expect the same $H$ to fit all galaxies/clusters.
This scenario still yields a falsifiable signature: the crossover radius extracted from kinematics
and the onset/extent of any lensing plateau should be consistent with a \emph{single} inferred
$H_{\rm eff}$ (or, in leaky extensions, with a pair $(H_{\rm eff},\ell_{\rm scr})$).

\paragraph{Crossover diagnostics.}
Operationally, $H$ (or $H_{\rm eff}$) can be inferred by any of:
(i) the slope-based knee $r_{\rm knee}\sim \mathcal{O}(H)$,
(ii) the rotation-curve plateau speed via $H=GM/v_\infty^2$, or
(iii) the lensing plateau height via $H=2\pi GM/(c^2\Delta\theta_\infty)$,
provided the relevant asymptotic regimes are reached.
Section~\ref{sec:leaky} emphasizes that imperfect confinement converts true asymptotes into finite
windows, so in practice one infers not only $H$ but also the \emph{persistence} of the slab regime.

\subsection{Joint inference from rotation curves and lensing}
\label{sec:cosmo_joint}

The optical identity \eqref{eq:lensing_force_identity} ties lensing directly to the same force law
that determines orbital dynamics.  In the sealed slab, the deep-regime predictions combine into the
parameter-free relation
\begin{equation}
  \Delta\theta_\infty = \frac{2\pi v_\infty^2}{c^2},
  \label{eq:cosmo_consistency}
\end{equation}
independent of $M$ and $H$ separately.  Equation~\eqref{eq:cosmo_consistency} is therefore the
cleanest ``mechanism test'' available in this framework: it checks whether the \emph{same} slab
regime supports both flattening and plateau lensing.

More generally, a joint analysis proceeds as follows:

\begin{enumerate}
\item \textbf{Kinematics:} fit $F(r)$ (or $v(r)$) to extract $r_{\rm knee}$ and an apparent plateau
      $v_\infty$ (or a finite quasi-flat window in leaky cases).

\item \textbf{Optics:} compute or measure $\Delta\theta(b)$ at large impact parameters, extract an
      apparent plateau $\Delta\theta_\infty$ (or a finite window/turnover).

\item \textbf{Consistency:} test \eqref{eq:cosmo_consistency} and check whether the inferred onset
      scales are compatible with a single $H_{\rm eff}$ (and, if leaky, with a single screening
      scale $\ell_{\rm scr}$).
\end{enumerate}

\paragraph{Leaky signature.}
In the leaky models of Section~\ref{sec:leaky}, the key observational handle is not only the plateau
height but its \emph{extent}.  For example, a Robin/lifted-zero-mode boundary predicts a plateau
only for $H\ll b\ll \ell_{\rm scr}$, followed by decay for $b\gtrsim \ell_{\rm scr}$.  In this way,
optics can constrain the confinement quality (via $k_0\simeq \ell_{\rm scr}^{-1}$ or an effective
reflectivity parameter) in a manner complementary to rotation curves.

\paragraph{Comment on scaling relations.}
The slab mechanism fixes the infrared functional form of the force, but the mapping between $M$ and
observables depends on how $H$ behaves across systems.
If $H$ is universal, then $v_\infty^2\propto M$ for isolated compact sources; if instead $H$ varies
systematically with environment, mass, or epoch, different scalings can emerge.  In either case,
the lensing--rotation consistency \eqref{eq:cosmo_consistency} remains the most direct test of the
slab regime itself, because it does not require independent mass calibration.

\subsection{Hydraulic Dark Energy as an optional module}
\label{sec:cosmo_de}

The slab confinement mechanism is a statement about the static scalar sector and does not, by itself,
fix the global expansion history.  However, the broader ``hydraulic cosmology'' interpretation of the
toy ontology suggests a natural way to discuss cosmic acceleration without introducing a fundamental
cosmological constant: treat the brane--bulk system as an \emph{open} hydraulic system with sinks
(throats) and sources (bulk injection).

At the level of phenomenological bookkeeping, let $V(t)$ denote the effective bulk volume associated
with the brane region under consideration and $\rho_0$ the preferred vacuum density of the stiff
superfluid.  If injection exceeds drainage, mass conservation implies
\begin{equation}
  \frac{\dd}{\dd t}\bigl(\rho_0 V\bigr) = \dot{Q}_{\rm in}-\dot{Q}_{\rm out}.
\end{equation}
If the vacuum remains approximately isodense/isobaric ($\rho_0\simeq \mathrm{const}$), then excess
injection manifests primarily as an increase in volume:
\begin{equation}
  \frac{\dot{V}}{V} \simeq \frac{\dot{Q}_{\rm in}-\dot{Q}_{\rm out}}{\rho_0 V},
\end{equation}
which can be interpreted as a macroscopic expansion rate of the brane-embedded domain.
In this picture, the constancy of the signal speed $c$ is tied to the constancy of the medium
properties (stiffness and density) rather than to a fundamental de Sitter vacuum.

We emphasize that this ``hydraulic dark energy'' module is conceptually separable from the slab dark
matter mechanism:
Sections~\ref{sec:greens}--\ref{sec:lensing} rely only on the static Poisson sector in a finite slab,
and the predictions \eqref{eq:dyn_vinf}, \eqref{eq:lensing_plateau}, and
\eqref{eq:cosmo_consistency} remain valid irrespective of the detailed expansion model.
A complete cosmological embedding would require specifying the source term(s) responsible for
$\dot{Q}_{\rm in}$, the equation of state maintaining $\rho_0$, and the coupling between injection,
bulk geometry, and the effective confinement scales $(H,\ell_{\rm scr})$.
Those model-building ingredients are deferred to future work; here we focus on the slab regime as a
minimal, falsifiable geometric mechanism for the apparent dark matter phenomenology.

\section{Predictions, constraints, and tests}
\label{sec:tests}

The slab confinement mechanism makes a small number of sharp, falsifiable statements.
Some are qualitative (the existence of a crossover), while others are quantitative and parameter
free once the system is in the deep slab regime.
This section collects the predictions in a form suitable for comparison with data and for guiding
future numerical/phenomenological work.

\subsection{Prediction box: sealed (Neumann) slab}
\label{sec:tests_prediction_box}

Assume the Poisson-sector scalar obeys Neumann confinement at $z=\pm H$ and that the system reaches
the deep slab regime at radii $r\gg H$ (and impact parameters $b\gg H$ for lensing).  Then:

\begin{enumerate}
\item \textbf{Force-law crossover at a geometric scale.}
The brane-plane acceleration transitions from Newtonian to slab scaling:
\begin{equation}
  F(r)\sim \frac{GM}{r^2}\quad (r\ll H),
  \qquad
  F(r)\sim \frac{GM}{H\,r}\quad (r\gg H),
\end{equation}
with a crossover (``knee'') radius $r_{\rm knee}=\mathcal{O}(H)$ identifiable by the log-slope
$\gamma(r)=\dd\ln F/\dd\ln r$ passing through $-3/2$.

\item \textbf{Flat rotation curves with normalization fixed by $H$.}
For circular motion,
\begin{equation}
  v^2(r)=rF(r)\;\to\; v_\infty^2=\frac{GM}{H}\qquad (r\gg H).
  \label{eq:tests_vinf}
\end{equation}

\item \textbf{Lensing plateau at large impact parameter.}
Using the optical mapping of Paper II, the bending angle becomes
\begin{equation}
  \Delta\theta(b)\to \Delta\theta_\infty = \frac{2\pi GM}{Hc^2},
  \qquad (b\gg H),
  \label{eq:tests_plateau}
\end{equation}
i.e.\ $\Delta\theta(b)$ saturates to an impact-parameter independent constant in the deep slab regime.

\item \textbf{Parameter-free lensing--rotation consistency.}
Combining \eqref{eq:tests_vinf} and \eqref{eq:tests_plateau} yields
\begin{equation}
  \boxed{\ \Delta\theta_\infty = \frac{2\pi v_\infty^2}{c^2}\ },
  \label{eq:tests_consistency}
\end{equation}
which is independent of $M$ and $H$ separately.  This is the cleanest direct test of the slab regime.

\item \textbf{Slope transitions as localized diagnostics.}
Define log-slopes
\begin{equation}
  \gamma(r)\equiv \frac{\dd\ln F}{\dd\ln r},
  \qquad
  \beta(b)\equiv \frac{\dd\ln \Delta\theta}{\dd\ln b}.
\end{equation}
Then
\begin{equation}
  \gamma:\ -2\to -1 \quad \text{across } r\sim H,
  \qquad
  \beta:\ -1\to 0 \quad \text{across } b\sim H,
\end{equation}
providing two independent, scale-localized crossover measurements.
\end{enumerate}

\subsection{Leakiness tests: when the slab regime is only a finite window}
\label{sec:tests_leaky}

If the bulk is not perfectly sealed, the Neumann zero mode is weakened or lifted.
Both leaky parameterizations in Section~\ref{sec:leaky} predict that the slab-like $1/r$ behavior is
not an asymptote but an intermediate window.
This is not a loophole; it is an additional set of falsifiable tests.

\paragraph{Model A (weighted images).}
With a reflectivity parameter $0<w<1$, the force is approximately slab-like only over a range
\begin{equation}
  H \ll r \ll r_{\max}(w)\sim 2n_{\rm eff}H \sim \frac{2H}{1-w},
\end{equation}
after which $F(r)$ departs from $1/r$.
Accordingly, rotation curves and lensing exhibit quasi-flat and quasi-plateau behavior only over a
finite dynamic range.  Measuring the extent of that range constrains $w$.

\paragraph{Model B (Robin / lifted zero mode).}
With a screening scale $\ell_{\rm scr}=k_0^{-1}$, slab-like behavior holds for
\begin{equation}
  H \ll r \ll \ell_{\rm scr},\qquad H \ll b \ll \ell_{\rm scr},
\end{equation}
and both the rotation-curve plateau and the lensing plateau are truncated for radii/impact
parameters beyond $\ell_{\rm scr}$.  Thus, joint dynamics+optics constrains both $H$ and
$\ell_{\rm scr}$ (equivalently $k_0$ or the Robin parameter $\alpha$).

\subsection{Constraints and practical inference strategy}
\label{sec:tests_inference}

A minimal data-analysis workflow suggested by the model is:

\begin{enumerate}
\item \textbf{Crossover scale from slopes.}
From kinematic data infer $r_{\rm knee}$ via $\gamma(r_{\rm knee})\approx -3/2$ and estimate
$H\sim r_{\rm knee}$ up to an order-unity calibration (which can be fixed by forward modeling with
the exact Green's function).

\item \textbf{Plateau normalization from dynamics and optics.}
If a clear kinematic plateau exists, infer $H$ from $H=GM/v_\infty^2$ (requiring an independent mass
estimate), or infer $H$ from $H=2\pi GM/(c^2\Delta\theta_\infty)$ if a lensing plateau exists.

\item \textbf{Mechanism test independent of mass calibration.}
If both a kinematic plateau and a lensing plateau are observed for the same system, test the
parameter-free relation \eqref{eq:tests_consistency} directly.

\item \textbf{Leakiness constraints from plateau extent.}
If plateaus are only approximate and terminate, extract the termination scale to constrain
$r_{\max}(w)$ (Model A) or $\ell_{\rm scr}$ (Model B).  A mismatch between the extent inferred from
dynamics and from optics would falsify the simplest leaky models.
\end{enumerate}

\subsection{What would falsify the slab confinement mechanism?}
\label{sec:tests_falsify}

The mechanism makes clear failure modes:

\begin{itemize}
\item \textbf{No consistent crossover scale:} if $\gamma(r)$ does not exhibit a robust transition
      from $-2$ toward $-1$ at any radius (after accounting for baryonic mass distribution and
      finite-size effects), the slab regime is not realized.

\item \textbf{Dynamics--optics inconsistency:} if systems with flat rotation curves show no tendency
      toward plateau lensing at large $b$, or if the observed plateaus violate
      \eqref{eq:tests_consistency} well beyond systematic uncertainties, the slab confinement picture
      (as implemented here) is disfavored.

\item \textbf{Incompatible plateau extents:} in leaky scenarios, if the radial extent of quasi-flat
      rotation and quasi-plateau lensing cannot be reconciled with a single confinement/leakiness
      scale, the simplest boundary-sensitivity models are ruled out.

\item \textbf{Strong-field contradictions:} if the inferred $H$ would force significant deviations
      from the already-validated near-field 1PN results of the earlier papers on scales where those
      results have been tested (e.g.\ Solar System), then $H$ cannot be a universal constant and must
      be an effective/environmental scale (or the slab mechanism is not applicable).
\end{itemize}

Taken together, these tests emphasize the intended status of the slab mechanism: it is a minimal,
geometric modification of the Poisson sector with a small number of parameters $(H)$ (or $(H,w)$ /
$(H,\ell_{\rm scr})$ in leaky extensions) and with sharply distinguishable optical signatures.

\section{Discussion and outlook}
\label{sec:discussion}

This paper isolates a simple but consequential geometric mechanism in the superfluid defect toy
model: solving the static scalar (Poisson) sector in a finite transverse slab with Neumann walls
modifies the global Green's function and induces an infrared crossover in the brane-plane force law.
The key results can be summarized compactly:

\begin{itemize}
\item \textbf{Dynamics:} a point source produces a crossover
\(
F(r)\sim GM/r^2
\)
for $r\ll H$ to
\(
F(r)\sim GM/(Hr)
\)
for $r\gg H$, implying a flat rotation-curve asymptote
\(
v_\infty^2=GM/H.
\)

\item \textbf{Optics:} using the same weak-field refractive-index mapping as in the optics paper of
the series, the slab regime predicts a large-impact-parameter bending plateau
\(
\Delta\theta_\infty = 2\pi GM/(Hc^2)
\),
and therefore a parameter-free lensing--rotation relation
\(
\Delta\theta_\infty=2\pi v_\infty^2/c^2.
\)

\item \textbf{Robustness/falsifiability:} the $1/r$ force law arises from the Neumann transverse
zero mode.  If confinement is imperfect, the zero mode is weakened or lifted, producing only a
finite intermediate $1/r$ window followed by eventual suppression.  This turns the boundary
condition from a hidden assumption into an observable diagnostic.
\end{itemize}

\subsection{What is proven here versus what is assumed}
The strongest part of the paper is the slab Green's function statement itself.
The force crossover and its normalization do not rely on phenomenological interpolation; they follow
from a well-posed boundary-value problem and are verified numerically in two independent
representations (images and KK/modes) with explicit convergence control.
Accordingly, the result should be viewed as a \emph{theorem about the scalar sector in a confined
geometry}, not as an empirical fitting ansatz.

At the same time, the cosmological interpretation of the slab thickness $H$ is necessarily an
assumption until the full brane--bulk dynamics that sets $H$ is specified.
This paper therefore treats $H$ operationally as the scale controlling the crossover and the plateau
normalizations, and it uses boundary-sensitivity models to parameterize imperfect confinement.
The resulting predictions remain falsifiable independent of whether $H$ is universal or
environment-dependent.

\subsection{Relation to other mechanisms in the toy model}
The slab mechanism modifies the infrared behavior without invoking additional nonbaryonic sources.
It should be distinguished from other routes to flattening that exist in the broader toy ontology,
notably velocity-dependent additive terms (the ``wake'' toy equation tested in the validation
script).  Those alternatives can also produce flat curves but generally differ in their optical
signatures, because lensing in the present framework is tied to the scalar potential gradients.
The lensing plateau derived here is therefore not a generic consequence of any flat rotation curve;
it is specific to a $F\propto 1/r$ slab regime and thus provides a meaningful discriminator.

\subsection{Immediate next steps}
The present paper was scoped to establish (i) the confined-geometry Green's function effect,
(ii) its numerical robustness, and (iii) its leading optical consequences.
The most direct extensions are:

\paragraph{(1) Complete the leaky-slab numerics and publish the diagnostic curves.}
Section~\ref{sec:leaky} provides two practical leakiness parameterizations.
The next concrete step is to generate a small set of ``robustness figures'':
log-slope $\gamma(r)$ curves for several $w$ or $k_0$, and the corresponding $\Delta\theta(b)$ curves
showing how the lensing plateau becomes a finite window.
This will convert the qualitative taxonomy into a quantitative constraint program.

\paragraph{(2) Extended sources and realistic baryonic profiles.}
While the point-source result cleanly isolates the mechanism, observations involve extended disks
and bulges.  Because the theory is linear in the Poisson sector, the extension is straightforward:
convolution of the surface density with the slab Green's function (or, in practice, numerical
integration using the mode representation).
A next paper (or an appendix expansion) should show how common baryonic profiles map into predicted
$F(r)$ and $\Delta\theta(b)$ curves, and how the knee scale compares to the observed onset of
flattening.

\paragraph{(3) Consistency across scales: galaxies to clusters.}
If $H$ is universal, the model makes strong across-system predictions; if $H$ is effective and
environment-dependent, one should articulate the physical control parameter(s) that set it
(e.g.\ epoch, ambient bulk pressure, brane curvature, throat density).
Either way, combining rotation and lensing across multiple mass scales will be the fastest route to
testing whether the slab regime is realized and whether leakage is required.

\paragraph{(4) Post-Newtonian and strong-field compatibility.}
The earlier 1PN papers in the series were designed to reproduce standard weak-field GR phenomenology
locally.  A key theoretical constraint is that the slab modification must not spoil those successes
on scales where they have already been validated.
This suggests that if $H$ is universal it must be sufficiently large to be irrelevant on Solar
System scales, or else the slab regime must be suppressed/absent locally (e.g.\ via effective
leakiness).
A systematic matching analysis between the confined Poisson sector and the established 1PN mapping
is therefore a natural follow-up.

\subsection{Broader outlook}
The main conceptual message is that apparent ``dark'' phenomenology can emerge from boundary
conditions and dimensional reduction rather than from new matter sources or ad hoc modifications of
the force law.
In the toy model, the confinement scale $H$ plays the role of a geometric parameter controlling the
infrared effective dimensionality of the scalar sector.
Because the optical sector couples to the same potential through $\ln N\simeq -2\Phi/c^2$, the model
predicts a tight link between rotation-curve flattening and large-impact-parameter lensing.
Whether nature realizes such a slab regime is an empirical question, but the predictions derived
here are sufficiently sharp that they can be ruled in or out with targeted joint dynamics+optics
comparisons.

Finally, the boundary sensitivity emphasized here should be viewed as a feature rather than a
liability: it provides a concrete handle on how ``sealed'' the bulk must be, and it suggests a path
toward interpreting deviations from idealized behavior (finite flat windows, truncated plateaus) as
measurements of confinement quality rather than as failures of the underlying mechanism.

\appendix

\section{Far-field normalization from a continuum approximation}
\label{app:farfield}

This appendix gives a short analytic derivation of the \emph{normalization} of the slab-regime force,
\begin{equation}
  F(r)\;\xrightarrow{\,r\gg H\,}\;\frac{\mu}{H\,r},
  \qquad \mu\equiv GM,
\end{equation}
directly from the image-charge representation.  The argument is a continuum approximation that
becomes accurate when many images contribute coherently, i.e.\ when the in-brane radius satisfies
$r\gg H$.  It is the simplest way to see why the coefficient is precisely $1/H$ rather than an
unspecified order-unity factor.

\subsection{Starting point: the Neumann image-stack force}
On the brane ($z=0$) the Neumann image construction gives
\begin{equation}
  F_{\rm img}(r)
  \;=\;
  \mu \sum_{n=-\infty}^{+\infty}
  \frac{r}{\left(r^2+(2nH)^2\right)^{3/2}}.
  \label{eq:app_farfield_Fimg}
\end{equation}
Define the smooth even kernel
\begin{equation}
  f(z)\equiv \frac{r}{(r^2+z^2)^{3/2}},
\end{equation}
so that \eqref{eq:app_farfield_Fimg} is
\begin{equation}
  F_{\rm img}(r) = \mu \sum_{n=-\infty}^{+\infty} f(2nH).
  \label{eq:app_farfield_sum}
\end{equation}

\subsection{Continuum approximation for $r\gg H$}
When $r\gg H$, the relevant values of $|z|$ that contribute substantially to $f(z)$ are of order
$|z|\lesssim r$.  Over that range the image spacing $2H$ is small compared to the scale on which
$f$ varies.  One may therefore approximate the discrete sum by a Riemann sum:
\begin{equation}
  \sum_{n=-\infty}^{+\infty} f(2nH)
  \;\approx\;
  \frac{1}{2H}\int_{-\infty}^{+\infty} f(z)\,\dd z
  \qquad (r\gg H).
  \label{eq:app_farfield_riemann}
\end{equation}
Substituting $f(z)$ and evaluating the integral gives
\begin{align}
  \int_{-\infty}^{+\infty} \frac{r}{(r^2+z^2)^{3/2}}\,\dd z
  &= \left[ \frac{z}{r\sqrt{r^2+z^2}} \right]_{-\infty}^{+\infty}
   = \frac{1}{r}-\left(-\frac{1}{r}\right)
   = \frac{2}{r}.
  \label{eq:app_farfield_integral}
\end{align}
Combining \eqref{eq:app_farfield_sum}--\eqref{eq:app_farfield_integral} yields the far-field force
with the correct coefficient:
\begin{equation}
  F_{\rm img}(r)
  \;\approx\;
  \mu \left(\frac{1}{2H}\right)\left(\frac{2}{r}\right)
  \;=\;
  \frac{\mu}{H\,r},
  \qquad (r\gg H).
  \label{eq:app_farfield_result}
\end{equation}
Equivalently,
\begin{equation}
  \boxed{\ Hr\,F(r)\to \mu \quad (r/H\to\infty)\ }.
\end{equation}
This is precisely the far-field normalization verified numerically in the main text.

\subsection{Corresponding potential: logarithm plus a gauge constant}
A similar continuum argument explains why the \emph{potential} in the slab regime is logarithmic
(up to an additive constant).
The brane-plane potential from the Neumann image stack can be written as
\begin{equation}
  \Phi(r,0) = -\mu \sum_{n=-\infty}^{+\infty}\frac{1}{\sqrt{r^2+(2nH)^2}} \;+\; \text{(constant)}.
  \label{eq:app_farfield_Phi_images}
\end{equation}
Replacing the sum by an integral as in \eqref{eq:app_farfield_riemann} gives
\begin{align}
  \sum_{n=-\infty}^{+\infty}\frac{1}{\sqrt{r^2+(2nH)^2}}
  &\approx \frac{1}{2H}\int_{-\infty}^{+\infty}\frac{\dd z}{\sqrt{r^2+z^2}}
   = \frac{1}{2H}\left[ 2\,\mathrm{asinh}\!\left(\frac{z}{r}\right)\right]_{-\infty}^{+\infty}.
\end{align}
The integral diverges logarithmically at the endpoints; this reflects the familiar fact that the
2D Green's function is defined only up to an additive constant.
Introducing a symmetric cutoff $|z|\le L$,
\begin{equation}
  \frac{1}{2H}\int_{-L}^{+L}\frac{\dd z}{\sqrt{r^2+z^2}}
  = \frac{1}{H}\,\mathrm{asinh}\!\left(\frac{L}{r}\right)
  = \frac{1}{H}\left[\ln\!\left(\frac{2L}{r}\right) + \mathcal{O}\!\left(\frac{r^2}{L^2}\right)\right].
\end{equation}
Absorbing the $L$-dependent piece into the arbitrary constant in \eqref{eq:app_farfield_Phi_images}
yields the slab-regime form
\begin{equation}
  \Phi(r,0)\;\simeq\; -\frac{\mu}{H}\,\ln r \;+\; \text{(constant)}.
  \label{eq:app_farfield_logPhi}
\end{equation}
Differentiating \eqref{eq:app_farfield_logPhi} recovers \eqref{eq:app_farfield_result}.

\subsection{Corrections and relation to the KK/mode form}
The continuum approximation captures the leading $m=0$ (zero-mode) contribution exactly.
A more precise statement is obtained by Poisson summation, which converts the discrete image stack
into the KK/mode expansion used in the main text:
\begin{equation}
  \sum_{n=-\infty}^{+\infty}
  \frac{r}{\left(r^2+(2nH)^2\right)^{3/2}}
  =
  \frac{1}{H\,r}
  + \frac{2\pi}{H^2}\sum_{m=1}^{\infty}
  m\,K_1\!\left(\frac{m\pi r}{H}\right).
\end{equation}
Since $K_1(x)$ decays exponentially for large $x$, the corrections to $1/(Hr)$ are
$\mathcal{O}(e^{-\pi r/H})$, consistent with the numerical behavior and with the interpretation
that the far field is controlled by the Neumann zero mode.

\section{Mode expansion details and zero-mode dominance}
\label{app:modes}

This appendix derives the KK/mode representation used in the main text and makes explicit why the
Neumann slab exhibits a $1/r$ far-field force on the brane.
The calculation is standard separation-of-variables for the Laplacian on a finite interval in $z$
combined with the 2D Green's functions in the brane plane.

\subsection{Neumann transverse eigenfunctions and normalization}

Let the slab domain be $z\in[-H,H]$ with Neumann walls.  The transverse eigenproblem is
\begin{equation}
  -f_m''(z)=k_m^2 f_m(z),\qquad \partial_z f_m(\pm H)=0,
\end{equation}
with orthonormality
\begin{equation}
  \int_{-H}^{H} f_m(z) f_n(z)\,\dd z = \delta_{mn}.
\end{equation}
A convenient orthonormal basis is
\begin{align}
  f_0(z) &= \frac{1}{\sqrt{2H}}, \qquad &&k_0=0,\\
  f_m(z) &= \frac{1}{\sqrt{H}}\cos\!\left(\frac{m\pi z}{H}\right), \qquad &&k_m=\frac{m\pi}{H},\quad m\ge 1.
\end{align}
For a brane source at $z=0$, only even modes appear; indeed $f_m(0)=1/\sqrt{H}$ for $m\ge 1$ and
$f_0(0)=1/\sqrt{2H}$.

The completeness relation yields an explicit series for the brane-localized delta function:
\begin{equation}
  \delta(z) = \sum_{m=0}^{\infty} f_m(0)\,f_m(z)
  = \frac{1}{2H} + \frac{1}{H}\sum_{m=1}^{\infty}\cos\!\left(\frac{m\pi z}{H}\right).
  \label{eq:app_delta_expansion}
\end{equation}

\subsection{Mode reduction of the slab Green's function}

Let $G_N(r,z)$ be the Neumann Green's function defined by
\begin{equation}
  \nabla^2 G_N(r,z) = 4\pi\,\delta^{(2)}(\mathbf{r})\,\delta(z),
  \qquad \partial_z G_N|_{z=\pm H}=0,
\end{equation}
where $r=\sqrt{x^2+y^2}$ is the brane-plane radius.
Expand
\begin{equation}
  G_N(r,z) = \sum_{m=0}^{\infty} g_m(r)\,f_m(z).
  \label{eq:app_G_expansion}
\end{equation}
Multiply by $f_n(z)$ and integrate over $z\in[-H,H]$. Using orthonormality and the eigenvalue
equation for $f_n$ gives the radial mode equations
\begin{equation}
  \left(\nabla_\perp^2 - k_n^2\right) g_n(r)
  = 4\pi\,\delta^{(2)}(\mathbf{r})\,f_n(0),
  \label{eq:app_radial_mode_eq}
\end{equation}
where $\nabla_\perp^2=\partial_x^2+\partial_y^2$ is the 2D Laplacian in the brane plane.

\subsection{2D Green's functions and explicit solutions}

We use the standard distributional identities in two dimensions:
\begin{align}
  \nabla_\perp^2 \ln r &= 2\pi\,\delta^{(2)}(\mathbf{r}),
  \label{eq:app_2d_log_identity}\\
  \left(\nabla_\perp^2 - k^2\right)K_0(k r) &= -2\pi\,\delta^{(2)}(\mathbf{r}),
  \qquad (k>0),
  \label{eq:app_2d_K0_identity}
\end{align}
where $K_\nu$ are modified Bessel functions of the second kind.

\paragraph{Zero mode ($m=0$).}
For $n=0$, $k_0=0$ and $f_0(0)=1/\sqrt{2H}$.  Equation \eqref{eq:app_radial_mode_eq} becomes
\begin{equation}
  \nabla_\perp^2 g_0(r) = \frac{4\pi}{\sqrt{2H}}\,\delta^{(2)}(\mathbf{r}).
\end{equation}
Using \eqref{eq:app_2d_log_identity}, one may take
\begin{equation}
  g_0(r) = \sqrt{\frac{2}{H}}\;\ln\!\left(\frac{r}{r_0}\right),
\end{equation}
where $r_0$ is an arbitrary reference scale reflecting the additive constant ambiguity of the 2D
Green's function.

\paragraph{Massive modes ($m\ge 1$).}
For $n=m\ge 1$, $k_m=m\pi/H$ and $f_m(0)=1/\sqrt{H}$.  Then \eqref{eq:app_radial_mode_eq} reads
\begin{equation}
  \left(\nabla_\perp^2 - k_m^2\right)g_m(r)=\frac{4\pi}{\sqrt{H}}\,\delta^{(2)}(\mathbf{r}).
\end{equation}
Using \eqref{eq:app_2d_K0_identity}, a convenient choice is
\begin{equation}
  g_m(r) = -\frac{2}{\sqrt{H}}\;K_0(k_m r),
  \qquad k_m=\frac{m\pi}{H}.
\end{equation}

\subsection{Green's function and potential on the brane}

Substituting the solutions into \eqref{eq:app_G_expansion} yields
\begin{equation}
  G_N(r,z)
  =
  \frac{1}{H}\ln\!\left(\frac{r}{r_0}\right)
  -\frac{2}{H}\sum_{m=1}^{\infty}
  K_0\!\left(\frac{m\pi r}{H}\right)\cos\!\left(\frac{m\pi z}{H}\right).
  \label{eq:app_G_closed}
\end{equation}
On the brane ($z=0$), $\cos(m\pi z/H)=1$, so
\begin{equation}
  G_N(r,0)
  =
  \frac{1}{H}\ln\!\left(\frac{r}{r_0}\right)
  -\frac{2}{H}\sum_{m=1}^{\infty}
  K_0\!\left(\frac{m\pi r}{H}\right).
  \label{eq:app_G_brane}
\end{equation}

The physical potential for a point source of strength $\mu\equiv GM$ is
\begin{equation}
  \Phi(r,z)=-\mu\,G_N(r,z),
\end{equation}
so on the brane
\begin{equation}
  \Phi(r,0)
  =
  -\mu\left[
    \frac{1}{H}\ln\!\left(\frac{r}{r_0}\right)
    -\frac{2}{H}\sum_{m=1}^{\infty}
    K_0\!\left(\frac{m\pi r}{H}\right)
  \right],
\end{equation}
which is equivalent (up to an additive constant absorbed into $r_0$) to the form quoted in the main
text.

Differentiating and using $\frac{\dd}{\dd x}K_0(x)=-K_1(x)$ gives the brane-plane inward acceleration
magnitude:
\begin{equation}
  F(r)\equiv -\partial_r\Phi(r,0)
  =
  \mu\left[
    \frac{1}{H\,r}
    +\frac{2\pi}{H^2}\sum_{m=1}^{\infty}
    m\,K_1\!\left(\frac{m\pi r}{H}\right)
  \right].
  \label{eq:app_F_modes}
\end{equation}

\subsection{Zero-mode dominance and exponential approach to the slab regime}

Equation \eqref{eq:app_F_modes} makes the slab physics explicit:
the Neumann zero mode contributes the term $1/(H r)$, while all higher KK modes are exponentially
suppressed in the far field because
\begin{equation}
  K_1(x) \sim \sqrt{\frac{\pi}{2x}}\,e^{-x}\left[1+\mathcal{O}\!\left(\frac{1}{x}\right)\right]
  \qquad (x\to\infty).
\end{equation}
Keeping only the leading $m=1$ correction for $r\gg H$ yields the useful estimate
\begin{equation}
  F(r)
  =
  \frac{\mu}{H\,r}\left[
    1
    + \mathcal{O}\!\left(\sqrt{\frac{r}{H}}\,e^{-\pi r/H}\right)
  \right],
  \label{eq:app_zero_mode_correction}
\end{equation}
so the approach to the slab asymptote is controlled by the first massive transverse gap
$\Delta k = \pi/H$.

\subsection{Near-field reconstruction of $1/r^2$}

For $r\ll H$, many KK modes contribute, and \eqref{eq:app_F_modes} reconstructs the local 3D Newtonian
behavior $F(r)\sim \mu/r^2$.
A compact way to see this is to use the exact identity relating the mode sum to the image stack
(derived via Poisson summation in the main text):
\begin{equation}
  \frac{1}{H\,r}
  +\frac{2\pi}{H^2}\sum_{m=1}^{\infty}
  m\,K_1\!\left(\frac{m\pi r}{H}\right)
  =
  \sum_{n=-\infty}^{+\infty}
  \frac{r}{\left(r^2+(2nH)^2\right)^{3/2}}.
  \label{eq:app_identity_repeat}
\end{equation}
When $r\ll H$, the $n=0$ image term dominates the right-hand side and gives $\mu/r^2$, with the
remaining images producing corrections suppressed by $(r/H)^3$ (as shown in
Section~\ref{sec:asymptotics_rotation}).  Thus, the KK representation inherits the correct near-field
normalization automatically while making the far-field zero-mode dominance manifest.

\section{Numerical methods and convergence parameters}
\label{app:numerics}

This appendix documents the numerical procedures used to evaluate the Neumann slab force law, to
control truncation errors in both representations (images and KK/modes), and to produce the
validation tables reported in Section~\ref{sec:numerics}.
All computations were performed in Mathematica using the accompanying script
\texttt{dark\_matter\_slab\_test\_revised.wl}.

\subsection{Core quantities and parameter conventions}

We evaluate the inward brane-plane acceleration magnitude for a point source of strength
$\mu\equiv GM$ in a slab of half-thickness $H$.
The two equivalent representations implemented are:

\paragraph{Image representation (Neumann walls).}
\begin{equation}
  F_{\rm img}(r)
  = \mu\sum_{n=-\infty}^{+\infty}
  \frac{r}{\left(r^2+(2nH)^2\right)^{3/2}}.
  \label{eq:app_num_Fimg}
\end{equation}

\paragraph{Mode/KK representation (Neumann spectrum).}
\begin{equation}
  F_{\rm mod}(r)
  = \mu\left[
    \frac{1}{H\,r}
    +\frac{2\pi}{H^2}\sum_{m=1}^{\infty}
    m\,K_1\!\left(\frac{m\pi r}{H}\right)
  \right].
  \label{eq:app_num_Fmod}
\end{equation}

Unless otherwise stated, the representative validation run uses
\begin{equation}
  H=10,\qquad \mu=1,
\end{equation}
and samples $r$ on a log-spaced grid over the interval $r\in[10^{-1},\,3\times 10^2]$, with additional
spot checks in the deep slab regime ($r\gg H$) to verify asymptotic normalization.

\subsection{Truncation and evaluation: images}

In practice we truncate the symmetric image sum at $|n|\le n_{\max}$:
\begin{equation}
  F_{\rm img}^{(n_{\max})}(r)
  = \mu\sum_{n=-n_{\max}}^{n_{\max}}
  \frac{r}{\left(r^2+(2nH)^2\right)^{3/2}}
  \;+\;\mu\,T_{\rm img}(r;n_{\max}),
  \label{eq:app_num_img_trunc}
\end{equation}
where $T_{\rm img}$ is an analytic tail correction approximating the contribution from $|n|>n_{\max}$.

\paragraph{Tail estimate.}
For large $|n|$, the summand admits an asymptotic expansion in powers of $(r/(2nH))^2$:
\begin{equation}
  \frac{r}{\left(r^2+(2nH)^2\right)^{3/2}}
  =
  \frac{r}{(2|n|H)^3}
  \left[
    1-\frac{3}{2}\left(\frac{r}{2|n|H}\right)^2
    +\mathcal{O}\!\left(\frac{r^4}{n^4H^4}\right)
  \right].
\end{equation}
Retaining the leading term yields the simple tail model
\begin{equation}
  T_{\rm img}(r;n_{\max})
  \approx 2\sum_{n=n_{\max}+1}^{\infty}\frac{r}{(2nH)^3}
  = \frac{r}{4H^3}\sum_{n=n_{\max}+1}^{\infty}\frac{1}{n^3}.
  \label{eq:app_num_img_tail}
\end{equation}
In the code, this tail is evaluated either via a direct high-precision numerical remainder
($\zeta(3)-H_{n_{\max}}^{(3)}$) or via its large-$n_{\max}$ approximation
$\sum_{n>n_{\max}}n^{-3}\sim 1/(2n_{\max}^2)$, depending on the requested speed/precision.
Higher-order terms in \eqref{eq:app_num_img_tail} can be included straightforwardly when extremely
large $r/H$ is probed and very tight relative error is desired.

\paragraph{Why tails matter most for $r\gg H$.}
In the deep slab regime, many images contribute coherently to build the $1/(Hr)$ asymptote, so the
raw truncated sum (without a tail) converges slowly as a function of $n_{\max}$.
The tail term in \eqref{eq:app_num_img_trunc} restores rapid convergence by capturing the net
contribution of the distant images analytically.

\subsection{Truncation and evaluation: modes}

We truncate the KK sum at $m\le m_{\max}$:
\begin{equation}
  F_{\rm mod}^{(m_{\max})}(r)
  = \mu\left[
    \frac{1}{H\,r}
    +\frac{2\pi}{H^2}\sum_{m=1}^{m_{\max}}
    m\,K_1\!\left(\frac{m\pi r}{H}\right)
  \right].
  \label{eq:app_num_mod_trunc}
\end{equation}

\paragraph{Exponential convergence for $r\gtrsim H$.}
Because $K_1(x)\sim \sqrt{\pi/(2x)}\,e^{-x}$ for large $x$, the truncation error for fixed $r$ is
dominated by the first omitted term and is approximately bounded by a geometric tail.
In practice this means:
\begin{itemize}
  \item for $r\gtrsim H$, only modest $m_{\max}$ is required (the zero mode plus a few corrections),
  \item for $r\ll H$, many modes contribute and $m_{\max}$ must be increased to reconstruct the local
        $1/r^2$ behavior to high precision.
\end{itemize}
This complementarity is one reason we keep both representations: the mode sum is extremely efficient
in the far field, while the image sum is conceptually transparent and efficient very near the source.

\subsection{Reference solutions and error metrics}

To produce convergence tables, we define ``reference'' forces against which truncated results are
compared:

\paragraph{Image reference.}
The image reference $F_{\rm ref,img}(r)$ is computed using a large cutoff $n_{\rm ref}$ and a
correspondingly tight tail evaluation.  The reported image truncation error is
\begin{equation}
  \varepsilon_{\rm img}(r;n_{\max})
  = \left|\frac{F_{\rm img}^{(n_{\max})}(r)-F_{\rm ref,img}(r)}{F_{\rm ref,img}(r)}\right|.
\end{equation}

\paragraph{Mode reference.}
The mode reference $F_{\rm ref,mod}(r)$ is computed with a large cutoff $m_{\rm ref}$ (and, when
needed, elevated working precision).  The reported mode truncation error is
\begin{equation}
  \varepsilon_{\rm mod}(r;m_{\max})
  = \left|\frac{F_{\rm mod}^{(m_{\max})}(r)-F_{\rm ref,mod}(r)}{F_{\rm ref,mod}(r)}\right|.
\end{equation}

\paragraph{Cross-check between representations.}
For each $r$ we also compute the representation disagreement
\begin{equation}
  \delta(r)
  = \left|\frac{F_{\rm mod}^{(m_{\max})}(r)-F_{\rm img}^{(n_{\max})}(r)}{F_{\rm img}^{(n_{\max})}(r)}\right|,
\end{equation}
using cutoffs chosen so that $\varepsilon_{\rm img}$ and $\varepsilon_{\rm mod}$ are each well below
the target tolerance (typically at or below numerical roundoff for the range of $r$ shown).

\subsection{Precision settings and numerical stability}

The computations are designed to be stable without excessive runtime:
\begin{itemize}
  \item For moderate $r/H$ the default machine-precision evaluation is sufficient, particularly for
        the mode sum.
  \item For deep slab checks at very large $r/H$ (e.g.\ verifying $HrF\to\mu$ at $r\sim 10^3$--$10^4$),
        we increase \texttt{WorkingPrecision} and use the mode sum as the primary evaluator, with the
        image sum serving as a cross-check using large $n_{\max}$ plus tail.
  \item All reported normalization checks use converged settings so that the displayed deviation from
        unity reflects physics (finite $r/H$) rather than truncation.
\end{itemize}

\subsection{Knee (crossover) estimation procedure}

To locate the transition scale without relying on plots, we compute the log-slope
\begin{equation}
  \gamma(r)=\frac{\dd\ln F}{\dd\ln r},
\end{equation}
using a symmetric finite difference on a log-spaced grid:
\begin{equation}
  \gamma(r_i)\approx
  \frac{\ln F(r_{i+1})-\ln F(r_{i-1})}{\ln r_{i+1}-\ln r_{i-1}}.
\end{equation}
We define the knee radius $r_{\rm knee}$ as the point where $\gamma$ crosses $-3/2$, obtained by
linear interpolation between the nearest grid points.
This definition is robust because $\gamma$ varies monotonically from $-2$ (Newtonian) to $-1$ (slab)
for the sealed Neumann case.

\subsection{Reproducibility notes}

The validation script is organized so that:
(i) $H$, $\mu$, $n_{\max}$, $m_{\max}$, and the sampled $r$-grid are set once at the top;
(ii) the same $r$ values are used for images, modes, and all error tables; and
(iii) the far-field and near-field normalization checks are printed explicitly in the form
$r^2F(r)\to\mu$ and $HrF(r)\to\mu$.
This makes it straightforward to repeat the convergence study for different $H$ values, to extend the
range of $r/H$, or to switch the boundary model (Section~\ref{sec:leaky}) while keeping the same
validation scaffolding.

\section{Lensing integral: stable forms and numerical checks}
\label{app:lensing}

This appendix collects numerically stable representations of the weak-field deflection integral used
in Section~\ref{sec:lensing}, and lists the analytic and internal-consistency checks we use to
validate the implementation.

\subsection{Starting point: deflection as a functional of the brane force}

Using the optical mapping summarized in Section~\ref{sec:inputs}, the deflection angle for a ray
with impact parameter $b$ can be written directly in terms of the inward acceleration magnitude
$F(R)$:
\begin{equation}
  \Delta\theta(b)
  \;\simeq\;
  \frac{2}{c^2}\int_{-\infty}^{+\infty}
  F\!\left(R(s)\right)\,\frac{b}{R(s)}\,\dd s,
  \qquad
  R(s)=\sqrt{b^2+s^2}.
  \label{eq:app_lens_base}
\end{equation}
All formulas below are equivalent rewritings of \eqref{eq:app_lens_base}.  The goal is to avoid
large cancellations and to keep the integrand smooth and bounded.

\subsection{Two stable coordinate transforms}

\paragraph{Form A (finite interval): $s=b\tan u$.}
Let $s=b\tan u$ so that $u\in(-\pi/2,\pi/2)$ and $R=b\sec u$.  Since $\dd s=b\sec^2u\,\dd u$ and
$b/R=\cos u$, Eq.~\eqref{eq:app_lens_base} becomes
\begin{equation}
  \Delta\theta(b)
  \;=\;
  \frac{2b}{c^2}\int_{-\pi/2}^{+\pi/2} F\!\left(b\sec u\right)\,\sec u\,\dd u.
  \label{eq:app_lens_u}
\end{equation}
Although $\sec u$ diverges at the endpoints, the \emph{full} integrand remains finite for all force
laws relevant here (Newtonian $F\sim 1/R^2$ and slab $F\sim 1/R$), so \eqref{eq:app_lens_u} is
numerically well behaved with standard adaptive quadrature.

\paragraph{Form B (no endpoint singularities): $s=b\sinh t$.}
An even smoother alternative is $s=b\sinh t$, for which $t\in(-\infty,\infty)$ and
$R=b\cosh t$.  Then $\dd s=b\cosh t\,\dd t$ and $b/R=1/\cosh t$, so the product $(b/R)\,\dd s=b\,\dd t$
and \eqref{eq:app_lens_base} reduces to the remarkably simple form
\begin{equation}
  \boxed{
  \Delta\theta(b)
  \;=\;
  \frac{2b}{c^2}\int_{-\infty}^{+\infty} F\!\left(b\cosh t\right)\,\dd t
  }
  \label{eq:app_lens_t}
\end{equation}
or, using evenness,
\begin{equation}
  \Delta\theta(b)
  \;=\;
  \frac{4b}{c^2}\int_{0}^{\infty} F\!\left(b\cosh t\right)\,\dd t.
  \label{eq:app_lens_t_even}
\end{equation}
This transform avoids divergent prefactors entirely; the integrand is smooth and decays rapidly in
$t$ for any force that falls at least as $1/R$.

In practice, we use \eqref{eq:app_lens_t} (or \eqref{eq:app_lens_t_even}) for production runs and
retain \eqref{eq:app_lens_u} as a cross-check.

\subsection{Analytic benchmarks (sanity checks)}

The lensing formulas admit two closed-form benchmarks that are also the asymptotes quoted in the
main text.

\paragraph{Newtonian benchmark ($F(R)=\mu/R^2$).}
Substituting $F(b\cosh t)=\mu/(b^2\cosh^2 t)$ into \eqref{eq:app_lens_t} gives
\begin{equation}
  \Delta\theta(b)
  =
  \frac{2b}{c^2}\int_{-\infty}^{+\infty}\frac{\mu}{b^2\cosh^2 t}\,\dd t
  =
  \frac{2\mu}{b\,c^2}\int_{-\infty}^{+\infty}\mathrm{sech}^2 t\,\dd t
  =
  \frac{4\mu}{b\,c^2},
\end{equation}
using $\int_{-\infty}^{\infty}\mathrm{sech}^2 t\,\dd t=2$.

\paragraph{Slab benchmark ($F(R)=\mu/(H R)$).}
Substituting $F(b\cosh t)=\mu/(H b\cosh t)$ into \eqref{eq:app_lens_t} yields
\begin{equation}
  \Delta\theta(b)
  =
  \frac{2b}{c^2}\int_{-\infty}^{+\infty}\frac{\mu}{H b\cosh t}\,\dd t
  =
  \frac{2\mu}{Hc^2}\int_{-\infty}^{+\infty}\mathrm{sech}\, t\,\dd t
  =
  \frac{2\pi\mu}{Hc^2},
\end{equation}
using $\int_{-\infty}^{\infty}\mathrm{sech}\,t\,\dd t=\pi$.

These two cases are used as first-pass unit tests: we implement $F(R)$ directly as $\mu/R^2$ and
$\mu/(HR)$ and verify that the numerical quadrature reproduces $4\mu/(b c^2)$ and $2\pi\mu/(Hc^2)$
over several decades in $b$.

\subsection{Numerical checks for the Neumann slab force}

For the physical slab force, we use either the converged KK/mode evaluator,
\begin{equation}
  F_{\rm mod}(r)
  = \mu\left[
    \frac{1}{H\,r}
    +\frac{2\pi}{H^2}\sum_{m=1}^{m_{\max}} m\,K_1\!\left(\frac{m\pi r}{H}\right)
  \right],
\end{equation}
or the converged image evaluator with tail correction (Appendix~\ref{app:numerics}).

We then perform three checks:

\paragraph{(1) Asymptotic recovery.}
Define the ratios
\begin{equation}
  \mathcal{R}_{\rm N}(b)\equiv \frac{\Delta\theta(b)}{4\mu/(b c^2)},
  \qquad
  \mathcal{R}_{\rm S}(b)\equiv \frac{\Delta\theta(b)}{2\pi\mu/(Hc^2)}.
\end{equation}
For a sealed Neumann slab, we should have $\mathcal{R}_{\rm N}(b)\to 1$ for $b\ll H$ and
$\mathcal{R}_{\rm S}(b)\to 1$ for $b\gg H$.  In practice we verify these limits over a wide dynamic
range with the same $(H,\mu)$ used in the force validation.

\paragraph{(2) Transform invariance.}
Compute $\Delta\theta(b)$ using both representations \eqref{eq:app_lens_u} and \eqref{eq:app_lens_t}
and verify agreement to the requested tolerance.  Because the transforms have different numerical
pathologies (bounded interval with endpoint structure versus unbounded interval with rapid decay),
agreement is a strong indicator that the evaluation is stable.

\paragraph{(3) Internal consistency with dynamics (plateau relation).}
In the deep slab regime, the model predicts
\begin{equation}
  \Delta\theta_\infty = \frac{2\pi v_\infty^2}{c^2},
\end{equation}
with $v_\infty^2=\mu/H$ extracted independently from the far-field force normalization (Section~\ref{sec:numerics}).
Numerically, we therefore compute $v_\infty^2$ from $rF(r)$ at large $r$ and compare
$2\pi v_\infty^2/c^2$ to the large-$b$ limit of $\Delta\theta(b)$.  This check is valuable because it
compares two independent integrals of the same underlying Green's function.

\subsection{Practical integration settings}

The following considerations make the evaluation reliable and reproducible:

\begin{itemize}
\item \textbf{Avoid $b=0$.}  For a true point mass, $F(R)$ diverges at $R=0$ and the thin-ray
      approximation is not meaningful at $b=0$.  All production runs use $b>0$.

\item \textbf{Prefer the $t$-transform for production.}
      Equation~\eqref{eq:app_lens_t} has no endpoint singularities and the integrand decays rapidly.
      One may truncate $t\in[-T,T]$ with $T$ chosen so that $b\cosh T$ lies well into the regime where
      $F(r)$ is already in its asymptotic form.

\item \textbf{Working precision.}
      When probing extreme ratios $b/H\gg 1$ or requiring tight relative errors on $\Delta\theta$,
      it is generally more efficient to compute $F(r)$ via the KK/mode representation (fast
      convergence in the far field) and to raise \texttt{WorkingPrecision} modestly for the
      quadrature.

\item \textbf{Consistency of cutoffs.}
      If $F(r)$ is evaluated from a truncated series, the $m_{\max}$ (or $n_{\max}$) choice should be
      made so that the force is converged at all radii $r=b\cosh t$ (or $r=b\sec u$) sampled by the
      integral.  This is automatic if the mode sum is used for $r\gtrsim H$ and the image method is
      reserved for near-field validation.
\end{itemize}

These checks and settings ensure that the lensing curves reported in the main text are properties of
the slab Green's function rather than artifacts of quadrature or truncation.

\section{Alternative flattening mechanism: wake term comparison}
\label{app:wake}

This appendix records a distinct flattening mechanism that arose in the toy-model exploration and was
validated numerically alongside the slab calculation.  It is \emph{not} part of the slab Green's
function result, but it is useful to keep on the record because it clarifies what aspects of the
phenomenology are specific to geometric confinement and which are not.

\subsection{Definition of the wake toy model (circular-orbit closure)}

The wake model is introduced at the level of the circular-orbit balance law by adding a
velocity-dependent term to the radial acceleration:
\begin{equation}
  \frac{v^2}{r}
  =
  \frac{\mu}{r^2}
  + C\,\frac{v}{r},
  \qquad \mu\equiv GM,
  \label{eq:wake_balance}
\end{equation}
where $C$ is a constant coupling with units of velocity.  Rearranging gives a quadratic equation
for the circular speed:
\begin{equation}
  v^2 - C v - \frac{\mu}{r}=0.
  \label{eq:wake_quadratic}
\end{equation}
The physical (positive) root is
\begin{equation}
  v(r)=\frac{1}{2}\left(C+\sqrt{C^2+\frac{4\mu}{r}}\right).
  \label{eq:wake_solution}
\end{equation}

\paragraph{Asymptotes.}
From \eqref{eq:wake_solution} one obtains:
\begin{align}
  r\to 0:\quad
  &v(r)=\sqrt{\frac{\mu}{r}}+\frac{C}{2}+\mathcal{O}\!\left(\sqrt{\frac{r}{\mu}}\,C^2\right),\\[4pt]
  r\to \infty:\quad
  &v(r)=C+\frac{\mu}{C\,r}+\mathcal{O}\!\left(\frac{1}{r^2}\right).
  \label{eq:wake_asymptotes}
\end{align}
Thus the model produces asymptotically flat rotation curves with
\begin{equation}
  v_\infty = C.
\end{equation}
The accompanying script verified this limit numerically (e.g.\ $v(r=10^8)\approx C$ for the tested
parameters).

\subsection{Conceptual distinction from slab confinement}

Although both the slab mechanism and the wake toy model can yield $v(r)\to\text{const}$, they differ
in three essential ways:

\paragraph{(1) Geometric versus dynamical origin.}
The slab flattening arises from a modified \emph{scalar Green's function} in a finite domain with
Neumann confinement; the far-field $F\propto 1/(Hr)$ law is a property of the boundary-value problem.
By contrast, the wake term in \eqref{eq:wake_balance} is \emph{explicitly velocity dependent} and is
not captured by a static scalar potential alone.

\paragraph{(2) Parameter interpretation and scaling.}
In the sealed slab,
\begin{equation}
  v_\infty^2=\frac{\mu}{H},
\end{equation}
so the plateau speed depends on the source strength $\mu$ and the geometric scale $H$.
In the wake toy model,
\begin{equation}
  v_\infty=C,
\end{equation}
independent of $\mu$ for sufficiently large $r$.
Accordingly, the two mechanisms imply different cross-system scalings unless $H$ (in the slab case)
or $C$ (in the wake case) is itself allowed to vary in a correlated way across environments.

\paragraph{(3) Crossover scale.}
In the wake model, a characteristic transition radius can be estimated by comparing the Newtonian
term to the wake term:
\begin{equation}
  \frac{\mu}{r^2}\sim C\frac{v}{r}.
\end{equation}
Using $v\sim C$ in the outer region gives
\begin{equation}
  r_{\rm tr}\sim \frac{\mu}{C^2},
\end{equation}
whereas in the slab case the transition is controlled geometrically by $r\sim H$.
This difference is observationally relevant: the slab knee is tied to a single spatial scale $H$,
while the wake knee scales with the source strength through $\mu/C^2$.

\subsection{Optical implications and why this appendix matters}

The main paper emphasized that slab confinement predicts a distinctive optical signature:
a large-impact-parameter deflection plateau
\begin{equation}
  \Delta\theta(b)\to \Delta\theta_\infty=\frac{2\pi \mu}{Hc^2},
\end{equation}
and therefore the parameter-free relation
\begin{equation}
  \Delta\theta_\infty=\frac{2\pi v_\infty^2}{c^2}.
\end{equation}
This is specific to a force law that truly behaves as $F(R)\propto 1/R$ in the regime relevant to the
lensing integral, and to the fact that in the optical sector of the series the deflection is tied to
the scalar potential via $\ln N\simeq -2\Phi/c^2$.

In contrast, the wake term in \eqref{eq:wake_balance} is not, as written, derivable from a static
scalar potential and therefore does not automatically inherit the slab lensing plateau prediction.
If the wake represents an effective ``lift'' force arising from a vector (gravitomagnetic-like)
sector, then its influence on photons requires a separate derivation from the \emph{full} emergent
metric (including any vector components) rather than from the scalar optical mapping alone.
Until that coupling is specified, the safe conclusion is:

\begin{quote}
Flat rotation curves by themselves do not uniquely identify slab confinement; the slab mechanism is
distinguished by its \emph{optical} prediction (the deflection plateau and the relation to
$v_\infty$).
\end{quote}

\subsection{Recommended paper positioning}

Because the wake toy model has a different theoretical status and (potentially) different optical
couplings, it should be presented separately from the slab Green's function claim.  Keeping it in an
appendix serves two purposes:
(i) it documents an alternative route to flattening that the toy model may support in other sectors,
and
(ii) it sharpens the interpretive point of the main paper: the slab mechanism earns credibility not
merely by producing flat curves, but by producing a \emph{linked} dynamics+optics signature that can
be tested (and falsified) in joint rotation+lensing analyses.


\end{document}
